\documentclass[11pt]{amsart}

\usepackage[letterpaper,margin=1.1in]{geometry}

\usepackage{amsmath,amssymb,mathtools,mathrsfs}
\usepackage{array}
\setcounter{MaxMatrixCols}{20}
\usepackage{microtype}
\usepackage{tikz}
\usepackage{float}
\usetikzlibrary{arrows.meta,calc,positioning,fit}
\usepackage[hidelinks]{hyperref}
\hypersetup{pdftitle={Cluster Algebras for Bosonic Plethysm},%
  pdfauthor={Yelin Fan and Jiarui Fei},%
  pdfsubject={Cluster structures, theta bases, and bosonic plethysm},%
  pdfkeywords={plethysm, cluster algebras, theta bases, invariant theory, polyhedral cones}}
\usepackage{quivers_plethysm}

\newtheorem{theorem}{Theorem}[section]
\newtheorem{proposition}[theorem]{Proposition}
\newtheorem{corollary}[theorem]{Corollary}
\newtheorem{lemma}[theorem]{Lemma}
\newtheorem*{theoremA}{Theorem A}
\newtheorem*{theoremB}{Theorem B}
\newtheorem*{theoremC}{Theorem C}
\theoremstyle{definition}
\newtheorem{definition}[theorem]{Definition}
\theoremstyle{remark}
\newtheorem{remark}[theorem]{Remark}
\newtheorem{example}[theorem]{Example}

\newcommand{\kk}{\Bbbk}
\newcommand{\GL}{\operatorname{GL}}
\newcommand{\Sym}{\operatorname{Sym}}
\newcommand{\SymMat}{\operatorname{SymMat}}
\newcommand{\htp}{\operatorname{ht}}
\newcommand{\cP}{\mathcal P}
\newcommand{\cR}{\mathcal R}
\newcommand{\cT}{\mathcal T}
\newcommand{\cU}{\mathcal U}
\newcommand{\frakS}{\mathfrak S}
\newcommand{\Mat}{\operatorname{Mat}}
\newcommand{\SL}{\operatorname{SL}}
\newcommand{\wt}{\operatorname{wt}}
\newcommand{\Frac}{\operatorname{Frac}}
\newcommand{\Rep}{\operatorname{Rep}}
\newcommand{\SI}{\operatorname{SI}}
\newcolumntype{L}[1]{>{\raggedright\arraybackslash}p{#1}}

\title{Cluster Algebras for Bosonic Plethysm}
\author{Yelin Fan}
\address{School of Mathematical Sciences, Shanghai Jiao Tong University, Shanghai, China}
\email{sjtu\_stanfan@sjtu.edu.cn}
\author{Jiarui Fei}
\address{School of Mathematical Sciences, Shanghai Jiao Tong University, Shanghai, China}
\email{jiarui@sjtu.edu.cn}
\subjclass[2020]{Primary 13F60; Secondary 05E10, 20G05, 52B20}
\keywords{plethysm, cluster algebras, theta bases, invariant theory, polyhedral cones}
\thanks{The authors were supported in part by the National Natural Science Foundation of China (Nos.~12131015 and 12571038).}
\thanks{An earlier version of this work was completed in 2022.}
\date{}

\begin{document}
\raggedbottom
\begin{abstract}
Let $\Bbbk$ be an algebraically closed field of characteristic zero, let
$V=\Bbbk^\ell$ and $W=\Bbbk^m$, and set
\[ \mathcal R_{\ell,m}=\operatorname{Sym}(\operatorname{Sym}^2V\otimes W)^{U_V}. \]
We construct an explicit skew-symmetrizable seed $\Sigma_{\ell,m}$ by
restricting and folding the determinantal seed for the flagged $m$-arrow
Kronecker quiver.  For every $\ell,m\ge2$, we have
\[ \mathcal R_{\ell,m}=\mathcal U(\Sigma_{\ell,m}), \]
with polynomial frozen coefficients, and $\Sigma_{\ell,m}$ admits a
reddening sequence.  The theta basis extends across the frozen boundary
exactly for parameters in a rational polyhedral cone $\mathscr C_{\ell,m}$.
Its weight fibers count the multigraded highest-weight multiplicities of
$\mathcal R_{\ell,m}$, and the Jacobi--Trudi identity expresses symmetric-square
plethysm coefficients as finite alternating sums of these counts.  
Optimized frozens give an explicit finite system of inequalities for $\mathscr C_{\ell,m}$.
\end{abstract}
\maketitle

\section{Introduction}
\label{sec:introduction}

Plethysm concerns the decomposition of a composite Schur functor
\[ \mathbf S_\mu(\mathbf S_\nu V) \cong \bigoplus_\lambda a_{\mu,\nu}^{\lambda}\,\mathbf S_\lambda V. \]
Here $\mu$ is the outer partition, $\nu$ is the inner partition, and
$\lambda$ is the resulting $\GL(V)$-highest weight.  The multiplicities
$a_{\mu,\nu}^{\lambda}$ are among the basic structural constants of polynomial
representation theory, but their general behavior remains poorly understood;
see, for example, \cite[Chapter~I]{Macdonald}.
This paper treats the case $\nu=(2)$.  Thus the inner functor is the
symmetric square, and we refer to this as the bosonic case.  The exterior
square gives the corresponding fermionic problem, studied in \cite{FFP}.

Fix $V=\kk^\ell$ and $W=\kk^m$; accordingly,
$\htp(\lambda)\leq\ell$ and $\htp(\mu)\leq m$.  The double highest-weight algebra
\[ \Sym(\Sym^2V\otimes W)^{U_V\times U_W} \]
contains the plethysm coefficients directly as joint weight
multiplicities.  It is useful, however, to retain the full $\GL(W)$-action
and consider instead
\begin{equation}
 \cR_{\ell,m}:=\Sym(\Sym^2V\otimes W)^{U_V}.
 \label{eq:intro-main-ring}
\end{equation}
After a basis of $W$ has been chosen, the algebra is graded by
$\alpha=(\alpha_1,\ldots,\alpha_m)\in\mathbb Z_{\ge0}^m$.  Its
$T_V$-weight space of weight $\lambda$ has dimension
\[ b_{\alpha,(2)}^\lambda =\dim\operatorname{Hom}_{\GL(V)}\!\left( \mathbf S_\lambda V, \bigotimes_{r=1}^m\Sym^{\alpha_r}(\Sym^2V) \right). \]
These multiplicities are nonnegative and retain the $\GL(W)$-symmetry that
is lost after taking $U_W$-invariants.  The Jacobi--Trudi identity recovers
plethysm from them:
\begin{equation}
 a_{\mu,(2)}^\lambda
 =\sum_{\sigma\in\frakS_m}\operatorname{sgn}(\sigma)
   b_{\alpha^\sigma(\mu),(2)}^\lambda,
 \qquad
 \alpha^\sigma(\mu)_i=\mu_i-i+\sigma(i).
 \label{eq:intro-JT}
\end{equation}
A term is understood to vanish when some component of
$\alpha^\sigma(\mu)$ is negative.  Our aim is therefore to describe the
weight spaces of \eqref{eq:intro-main-ring} by cluster-theoretic means.

\subsection{Main results}

Write $X_{\ell,m}=\SymMat_\ell^m$.  The algebra
$\cR_{\ell,m}$ is the invariant ring $\kk[X_{\ell,m}]^{U_\ell}$ for the
simultaneous congruence action
\[ u\cdot(C_1,\ldots,C_m) =(u^{-T}C_1u^{-1},\ldots,u^{-T}C_mu^{-1}). \]
For each $\ell,m\ge2$ we define a seed $\Sigma_{\ell,m}$ whose extended
cluster consists of the determinants $\Delta_r=\det C_r$ and a family of
chamber determinants $z_{i,j}^{(n)}$.  Its frozen set has $m+\ell-1$
vertices, and its mutable set has
\[ \frac{(m-1)(\ell-1)(\ell+2)}2 \]
vertices.  The exchange relations are given explicitly in
Theorem~\ref{thm:folded-relations}.
The seed $\Sigma_{4,3}$ is shown in Figure~\ref{fig:seed-43}.

\begin{theoremA}
For every $\ell,m\ge2$, the seed $\Sigma_{\ell,m}$ has full row rank and
skew-symmetrizable principal part, and
\[ \cR_{\ell,m}=\mathcal U(\Sigma_{\ell,m}) \]
with polynomial frozen coefficients.
\end{theoremA}

\begin{theoremB}
For every $\ell,m\ge2$, the seed $\Sigma_{\ell,m}$ admits a reddening
sequence.
\end{theoremB}

Theorem~B is proved in Theorem~\ref{thm:reddening}.  For odd $m$, the
mutable exchange matrix is reduced, by mutations and source--sink
deletions, to a string-diagram matrix of type $A_{\ell-1}$.  The even case
then follows because $B_{\ell,m}^{\mathrm{uf}}$ is a principal submatrix of
$B_{\ell,m+1}^{\mathrm{uf}}$.

Let $M_{\ell,m}$ be the character lattice of the initial cluster torus, and
let $D_f$ be the prime divisor associated with a frozen variable $z_f$.
For a theta function $\vartheta_g$, put
\[ \mathscr C_{\ell,m} =\{g\in M_{\ell,m,\mathbb R}: \operatorname{ord}_{D_f}(\vartheta_g)\ge0 \text{ for every frozen index }f\}. \]
The boundary valuations are integral convex piecewise-linear functions, so
$\mathscr C_{\ell,m}$ is a rational polyhedral cone.  Let
$\mathsf W_{\ell,m}$ denote the weight map of the initial seed and set
\[ \mathscr P_{\ell,m}(\lambda;\alpha) =\{g\in\mathscr C_{\ell,m}: \mathsf W_{\ell,m}g=(\lambda;\alpha)\}. \]

\begin{theoremC}
For every $\ell,m\ge2$, the family
\[ \{\vartheta_g:g\in\mathscr C_{\ell,m}\cap M_{\ell,m}\} \]
is a basis of $\cR_{\ell,m}$.  Moreover,
\[ b_{\alpha,(2)}^\lambda =\bigl|\mathscr P_{\ell,m}(\lambda;\alpha)\cap M_{\ell,m}\bigr|, \]
and hence
\[ a_{\mu,(2)}^\lambda =\sum_{\sigma\in\frakS_m}\operatorname{sgn}(\sigma) \bigl|\mathscr P_{\ell,m} (\lambda;\alpha^\sigma(\mu))\cap M_{\ell,m}\bigr|. \]
\end{theoremC}

The basis statement is Theorem~\ref{thm:theta-basis}, and the two formulas
are Corollary~\ref{cor:polyhedral-b} and
Theorem~\ref{thm:polyhedral-plethysm}.  Thus the coefficients
$b_{\alpha,(2)}^\lambda$ have a lattice-point interpretation, while plethysm is
obtained from these counts by the signed Jacobi--Trudi formula
\eqref{eq:intro-JT}.  We do not obtain a positive rule for
$a_{\mu,(2)}^\lambda$ itself.

When $\ell$ is odd and $m$ is even, every frozen boundary divisor has an
optimized seed.  Pulling the corresponding chiral-dual monomial back to the
initial chart gives a finite Laurent polynomial, and its tropicalization
gives the boundary inequality.  This produces an integral matrix
$H_{\ell,m}$ with
\[ \mathscr C_{\ell,m} =\{g:H_{\ell,m}g\ge0\}; \]
see Proposition~\ref{prop:optimized-boundary-seeds} and
Corollary~\ref{cor:boundary-cone-optimized-seeds}.  The cone and theta basis
exist for all parities, but the uniform optimized-seed construction is only
proved in this parity range.

Proposition~\ref{prop:successive-matrix-face} gives compatible face
models for different numbers of matrices.  For $2\le r<m$, the
simultaneous zero-degree locus for $C_{r+1},\ldots,C_m$ is an exposed face
of $\mathscr C_{\ell,m}$, and the theta functions indexed by its integral
points form a basis of $\cR_{\ell,r}$.  In particular, when $\ell$ is
odd, the explicit inequalities for $2n$ matrices restrict to a polyhedral
model for $2n-1$ matrices.

\subsection{The folded determinantal seed}

The construction begins with the two-sided matrix invariant ring
\[ \cT_{\ell,m} =\kk[\Mat_\ell^m]^{U_\ell^{\mathrm L}\times U_\ell^{\mathrm R}}. \]
The cluster structure on the semi-invariant ring of the flagged Kronecker
quiver in \cite{FeiKroneckerII} transfers to $\cT_{\ell,m}$.  Its cluster
functions are determinants of matrices formed from alternating products of
the $A_r$ and $A_r^{-1}$; determinant factors clear the apparent
denominators, so these are polynomial invariants.  Matrix transposition
interchanges the two flag arms and acts on this seed.

Restriction from arbitrary matrices to symmetric matrices sends the left
and right chamber functions in each transpose pair to the same determinant.
This identification of functions does not by itself determine an exchange
matrix.  One must prove that transposition is admissible and then sum the
columns over each transpose orbit.  Fixed chambers have orbit size one and
the remaining chambers have orbit size two; their interaction produces the
valuations $(1,2)$ and $(2,1)$ in the folded quiver.  Consecutive triangular
blocks are identified along alternating sides: for odd $n$ the side
$i+j=\ell$ is shared with block $n-1$, whereas for even $n$ the shared side
is $i=0$.  The orbit-sum construction gives a skew-symmetrizer equal to the
orbit-size function.

Algebraic independence is proved on a Gauss open set.  The fixed torus of
the two-sided cluster chart is a dense open subset of
$T_\ell\times\SymMat_\ell^{m-1}$, and its coordinate functions are precisely
the folded initial family.  To prove the upper-cluster equality, the first
inclusion is obtained from the maximal-rank upper-bound theorem.  The
invariant ring is factorial; irreducibility of the initial functions and
coprimality with their one-step mutations imply
$\mathcal U(\Sigma_{\ell,m})\subseteq\cR_{\ell,m}$.  For the reverse
inclusion, the $LDL^T$ decomposition gives generators after localization at
the leading principal minors of $C_1$.  Restriction from the two-sided
matrix algebra is surjective on this open set, and the same argument applies
to $C_2$.  Intersecting the two localizations removes all denominators and
yields Theorem~A.  In particular, the proof does not require permanent
inversion of the central determinants.

The reddening argument uses the disk quivers and flip sequences of Fock and
Goncharov \cite{FockGoncharov}.  A transpose-fixed half-diamond diagonal is
removed directly in the folded matrix.  Between two such removals, the two
flip regions are disjoint and exchanged by transposition, so their mutation
sequences descend by orbit mutation.  This reduces the odd-$m$ case to a
string diagram; the result for even $m$ follows by passage to a principal
submatrix.

Finally, the reddening sequence and full row rank imply the full
Fock--Goncharov condition in the sense of
Gross--Hacking--Keel--Kontsevich \cite{GHKK}.  The open cluster variety
therefore has a theta basis indexed by all integral tropical points.  The
polynomial frozen coefficients define a partial compactification whose
boundary consists of the frozen prime divisors.  Valuative independence and
theta reciprocity \cite{CMMM} select exactly the theta functions regular
along this boundary, giving Theorem~C.

\subsection{Relation to earlier work and provenance of the results}

Cluster algebras and upper cluster algebras were introduced in
\cite{FZ1,BFZ}; the coefficient formalism used here is compatible with
\cite{FZ4}.  Cluster structures on semi-invariant rings of triple flags and
flagged Kronecker quivers were constructed in
\cite{Fsemi,FeiKroneckerII}.  The sink--source extension and frozen
amalgamation framework of \cite{FeiWeymanExtension} gives a complementary
construction of the simply-laced hive quivers for complete $m$-tuple flags.
It describes the triangular pieces underlying the unfolded seed, but it does
not imply the symmetric restriction or the transpose-orbit fold considered
here.  Admissible orbit folding is part of the standard theory of
non-simply-laced cluster algebras; see \cite{DupontFolding}.  Recent work of
Ye constructs folded cluster structures on $SL_n/SO_n$ and on the affine
space of a single symmetric matrix \cite{YeSymmetric}.  Those varieties and
seeds are distinct from the highest-weight invariant algebras of tuples of
symmetric matrices considered here.

The representation-theoretic identities in
Section~\ref{sec:plethysm-highest-weight} are standard consequences of the
Cauchy decomposition, highest-weight theory, and the Jacobi--Trudi identity;
proofs are included to fix the gradings and sign conventions.  The two-sided
transfer and the unfolded seed come from
\cite{Grosshans,FeiKroneckerII}, the upper-bound argument uses \cite{GSV},
the reddening argument uses \cite{FockGoncharov,CaoString,CaoLi}, and the
theta-basis argument uses \cite{GHKK,CMMM}.

To the authors' knowledge, the following results are new: the symmetric
restriction and explicit orbit-sum folded seed for all $\ell,m$; the folded
exchange relations, full-rank statement, and equality with
$\cR_{\ell,m}$; the reddening construction for every $\ell,m\ge2$; the
boundary theta basis and its plethysm fibers; and the optimized boundary
seeds in the stated parity range.  None of these conclusions is a formal
specialization of the flagged Kronecker construction: transpose-related
functions coalesce, fixed orbits create valued arrows, and both the
exchange matrix and the upper-cluster equality require separate arguments.

\subsection{Organization}

Section~\ref{sec:plethysm-highest-weight} relates plethysm to the
multigraded highest-weight multiplicities of $\cR_{\ell,m}$.
Section~\ref{sec:gaussian-localization} describes the symmetric Gauss chart
and its polynomial coordinate lifts.  Section~\ref{sec:matrix-invariants}
transfers the flagged Kronecker seed to two-sided matrix invariants and
restricts the cluster functions to symmetric matrices.
Section~\ref{sec:folding} constructs the folded exchange matrix and proves
the exchange relations.  Section~\ref{sec:upper-cluster-equality} proves
Theorem~A, Section~\ref{sec:reddening} proves Theorem~B, and
Section~\ref{sec:polyhedral-plethysm} proves Theorem~C, the optimized
boundary results, and the degree-zero face construction.  The appendices
supply the transfer dictionary, the remaining coprimality arguments, the
folded half-diamond calculation, and the even boundary-path mutation.

\section{Plethysm and highest-weight invariants}
\label{sec:plethysm-highest-weight}

The identities in this section are standard consequences of the Cauchy
formula, highest-weight theory, and the Jacobi--Trudi identity; see
\cite[Chapter~I]{Macdonald}.  We include the proofs in order to fix the
multigrading and sign conventions used later.

Throughout the paper, $\kk$ is an algebraically closed field of
characteristic zero.  If $E$ is a finite-dimensional $\kk$-vector space, we
write
\[ G_E=\GL(E). \]
After choosing an ordered basis of $E$, let $B_E=T_EU_E$ be the standard
upper-triangular Borel subgroup, with diagonal torus $T_E$ and upper
unitriangular subgroup $U_E$.

A partition is a weakly decreasing sequence of nonnegative integers with
finite support.  For a partition $\rho$, write $|\rho|$ for its size and
$\htp(\rho)$ for the number of its nonzero parts.  In plethysm, $\mu$ is the
outer partition, $\nu$ the inner partition, and $\lambda$ the resulting
$G_V$-highest weight.  The symbol
$\alpha=(\alpha_1,\ldots,\alpha_m)\in\mathbb Z_{\geq0}^m$ denotes a weak
composition.  We write $\mathbf S_\rho$ for the Schur functor associated
with $\rho$; in particular, $\mathbf S_{(d)}E=\Sym^d(E)$.

Let $V$ and $W$ be vector spaces of dimensions $\ell$ and $m$, respectively,
and let $\nu$ be a partition with $\htp(\nu)\leq \ell$.  We consider the
polynomial algebra
\begin{equation}
 \cP_\nu(V,W)
 :=\Sym\bigl(\mathbf S_\nu V\otimes W\bigr).
 \label{eq:ambient-polynomial-algebra}
\end{equation}
Equivalently,
\[ \cP_\nu(V,W) =\kk\bigl[\mathbf S_\nu(V^\vee)\otimes W^\vee\bigr], \]
so \eqref{eq:ambient-polynomial-algebra} may be viewed either as a symmetric
algebra or as the coordinate ring of the dual representation.  The group
$G_V\times G_W$ acts naturally on this algebra.  Its $U_V$-invariant
subalgebra
\[ \cR_\nu(V,W):=\cP_\nu(V,W)^{U_V} \]
will be called the \emph{highest-weight algebra} associated with
$\mathbf S_\nu V\otimes W$.

\subsection{Plethysm coefficients as double highest-weight multiplicities}

For partitions $\lambda$ and $\mu$, with
$\htp(\lambda)\leq\ell$ and $\htp(\mu)\leq m$, the plethysm coefficient
$a_{\mu,\nu}^{\lambda}$ is defined by
\[ \mathbf S_\mu(\mathbf S_\nu V) \cong \bigoplus_{\rho} a_{\mu,\nu}^{\rho}\,\mathbf S_\rho V. \]
Under the stated height assumptions, this multiplicity is independent of the
choice of $V$.

For a $T_V\times T_W$-module $M$ and dominant weights $\lambda$ of $G_V$ and
$\mu$ of $G_W$, let $M_{\lambda;\mu}$ denote the corresponding joint
weight space.

The Cauchy formula and highest-weight theory give
\[ a_{\mu,\nu}^{\lambda}=\dim_\kk\left(\cP_\nu(V,W)^{U_V\times U_W}\right)_{\lambda;\mu}; \]
see \cite[Chapter~I]{Macdonald}.

The double invariant algebra records plethysm coefficients directly.  We
instead retain the full $G_W$-action and study $\cR_\nu(V,W)$.

\subsection{Multigraded multiplicities}

Fix a basis $w_1,\ldots,w_m$ of $W$.  The resulting decomposition
$W=\bigoplus_{r=1}^m\kk w_r$ gives $\cP_\nu(V,W)$ and
$\cR_\nu(V,W)$ a $\mathbb Z_{\geq0}^m$-grading.  For
$\alpha=(\alpha_1,\ldots,\alpha_m)$, the corresponding homogeneous component
is naturally isomorphic to
\begin{equation}
 \cP_\nu(V,W)_\alpha
 \cong
 \bigotimes_{r=1}^m
 \Sym^{\alpha_r}(\mathbf S_\nu V).
 \label{eq:multidegree-piece}
\end{equation}
For a partition $\lambda$ with $\htp(\lambda)\leq\ell$, define
\begin{equation}
 b_{\alpha,\nu}^{\lambda}
 :=
 \dim_\kk
 \left(\cR_\nu(V,W)_\alpha\right)_\lambda.
 \label{eq:b-definition}
\end{equation}
We refer to these numbers as the \emph{multigraded highest-weight
multiplicities}.  By convention, $b_{\alpha,\nu}^{\lambda}=0$ if any component of
$\alpha$ is negative.

By \eqref{eq:multidegree-piece} and highest-weight theory,
\begin{equation}
 b_{\alpha,\nu}^{\lambda}
 =\dim_\kk\operatorname{Hom}_{G_V}\!\left(\mathbf S_\lambda V,
   \bigotimes_{r=1}^m\Sym^{\alpha_r}(\mathbf S_\nu V)\right).
 \label{eq:b-as-multiplicity}
\end{equation}

\begin{remark}
Although $\alpha$ is retained as an ordered $T_W$-weight, the multiplicity
$b_{\alpha,\nu}^{\lambda}$ is unchanged when the entries of $\alpha$ are
permuted.
\end{remark}

\subsection{Recovering plethysm coefficients}

The multiplicities \eqref{eq:b-definition} do not equal plethysm coefficients
in general.  They nevertheless determine them by a finite alternating sum.
For a partition $\mu$ with $\htp(\mu)\leq m$, pad $\mu$ with
zeros and write
$\mu=(\mu_1,\ldots,\mu_m).$
For $\sigma\in\frakS_m$, define an integral $m$-tuple
\begin{equation}
 \alpha^\sigma(\mu)_i
 :=\mu_i-i+\sigma(i),
 \qquad 1\leq i\leq m.
 \label{eq:shifted-composition}
\end{equation}
The equality $|\alpha^\sigma(\mu)|=|\mu|$ holds, but
$\alpha^\sigma(\mu)$ need not be a weak composition.

\begin{theorem}
\label{thm:alternating-plethysm}
For partitions $\lambda,\mu,\nu$ satisfying
$\htp(\lambda)\leq\ell$, $\htp(\mu)\leq m$, and
$\htp(\nu)\leq\ell$, one has
\begin{equation}
 a_{\mu,\nu}^{\lambda}
 =
 \sum_{\sigma\in\frakS_m}
 \operatorname{sgn}(\sigma)\,
 b_{\alpha^\sigma(\mu),\nu}^{\lambda}.
 \label{eq:alternating-plethysm}
\end{equation}
Terms for which \eqref{eq:shifted-composition} has a negative component are
understood to be zero.
\end{theorem}

\begin{proof}
Apply the plethystic ring homomorphism $f\mapsto f[s_\nu]$ to the
Jacobi--Trudi identity
\[ s_\mu=\det\bigl(h_{\mu_i-i+j}\bigr)_{1\le i,j\le m} \]
and expand the determinant.  By \eqref{eq:b-as-multiplicity},
the coefficient of $s_\lambda$ in the term indexed by $\sigma$ is
$b_{\alpha^\sigma(\mu),\nu}^{\lambda}$.  Comparing coefficients proves
\eqref{eq:alternating-plethysm}.
\end{proof}

\begin{example}
When $m=2$ and $\mu=(\mu_1,\mu_2)$, Theorem
\ref{thm:alternating-plethysm} reads
\[ a_{\mu,\nu}^{\lambda} = b_{(\mu_1,\mu_2),\nu}^{\lambda} - b_{(\mu_1+1,\mu_2-1),\nu}^{\lambda}. \]
\end{example}

\subsection{The symmetric-square specialization}

We call plethysm with inner partition $(2)$, equivalently with the symmetric-square
functor $\Sym^2$, the \emph{bosonic} case.  The exterior-square analogue,
which gives fermionic plethysm, is treated in \cite{FFP}.  For the remainder
of the paper we specialize to $\nu=(2)$ and set
\[ V=\kk^\ell, \qquad W=\kk^m. \]
We abbreviate
\begin{equation}
 \cR_{\ell,m}
 :=\cR_{(2)}(V,W)
 =\Sym\bigl(\Sym^2V\otimes W\bigr)^{U_V}.
 \label{eq:main-ring-abstract}
\end{equation}
To obtain the matrix realization, let
\[ \SymMat_\ell :=\{C\in\operatorname{Mat}_{\ell\times\ell}(\kk):C^T=C\}, \qquad X_{\ell,m}:=\SymMat_\ell^m. \]
After choosing the standard bases of $V$ and $W$, the algebra in
\eqref{eq:main-ring-abstract} identifies with
\begin{equation}
 \cR_{\ell,m}=\kk[X_{\ell,m}]^{U_\ell},
 \label{eq:main-ring-matrices}
\end{equation}
where $U_\ell$ is the upper unitriangular subgroup of $\GL_\ell$ and
\[ u\cdot(C_1,\ldots,C_m) = \bigl(u^{-T}C_1u^{-1},\ldots,u^{-T}C_mu^{-1}\bigr), \qquad u^{-T}:=(u^{-1})^T. \]
For $0\le i\le\ell$, set
\begin{equation}
 h_i:=\det C_1[1,\ldots,i\mid1,\ldots,i],
 \qquad h_0:=1,
 \label{eq:hi-definition}
\end{equation}
and put
\[ H_1:=\prod_{i=1}^{\ell}h_i. \]
Theorem~\ref{thm:alternating-plethysm} recovers plethysm coefficients for
inner partition $(2)$ from the multigraded weight multiplicities of
\eqref{eq:main-ring-matrices}.





\section{Gauss-decomposition coordinates on the invariant ring}
\label{sec:gaussian-localization}

On the principal open set defined by the leading principal minors, the
$U_\ell$-action has a cross-section described by the $LDL^T$ decomposition.

\subsection{The Gauss chart}
\label{subsec:gaussian-slice}

Let
\[ X_{\ell,m}^{\mathrm G} :=\{(C_1,\ldots,C_m)\in X_{\ell,m}:H_1(C)\ne0\}. \]
Every $C_1$ occurring in this open set has a unique $LDL^T$
decomposition
\begin{equation}
 C_1=u^TDu,
 \qquad u\in U_\ell,
 \qquad D\in T_\ell.
 \label{eq:symmetric-gaussian-decomposition}
\end{equation}
For $C=(C_1,\ldots,C_m)\in X_{\ell,m}^{\mathrm G}$, set
\[ S_r:=u^{-T}C_ru^{-1}\in\SymMat_\ell, \qquad 2\le r\le m. \]

\begin{proposition}[Gauss chart]
\label{prop:gaussian-slice}
The morphism
\begin{align}
 \Psi:T_\ell\times\SymMat_\ell^{m-1}\times U_\ell
 &\longrightarrow X_{\ell,m}^{\mathrm G},
 \notag\\
 (D,S_2,\ldots,S_m,u)
 &\longmapsto
 (u^TDu,u^TS_2u,\ldots,u^TS_mu)
 \label{eq:gaussian-slice-map}
\end{align}
is an isomorphism.  Under this isomorphism the congruence action of
$w\in U_\ell$ is right translation on the last factor:
\begin{equation}
 w\cdot(D,S_2,\ldots,S_m,u)
 =(D,S_2,\ldots,S_m,uw^{-1}).
 \label{eq:gaussian-action}
\end{equation}
Consequently,
\begin{equation}
 \cR_{\ell,m}[H_1^{-1}]
 \cong\kk[T_\ell\times\SymMat_\ell^{m-1}].
 \label{eq:gaussian-quotient-ring}
\end{equation}
\end{proposition}

\begin{proof}
The decomposition \eqref{eq:symmetric-gaussian-decomposition} is the usual
$LDL^T$ decomposition, written with the upper unitriangular factor
$u=L^T$.  Its diagonal entries are
\begin{equation}
 D_{aa}=\frac{h_a}{h_{a-1}},
 \qquad 1\le a\le\ell.
 \label{eq:gaussian-diagonal-pivots}
\end{equation}
Gaussian elimination expresses $u$ regularly on the principal open
$H_1\ne0$, so the assignment
\[ C\longmapsto(D,S_2,\ldots,S_m,u) \]
is the inverse of \eqref{eq:gaussian-slice-map}.  Formula
\eqref{eq:gaussian-action} follows from the congruence action.  Taking
invariants removes the $U_\ell$-factor, and localization commutes with
invariants because $H_1$ is invariant.  This proves
\eqref{eq:gaussian-quotient-ring}.
\end{proof}

For $2\le r\le m$ and $1\le a\le b\le\ell$, let
\[ \xi_{r;a,b}(D,S_2,\ldots,S_m):=(S_r)_{ab}. \]
Since the $h_i$ are Laurent coordinates on $T_\ell$ by
\eqref{eq:gaussian-diagonal-pivots}, Proposition~\ref{prop:gaussian-slice}
gives
\begin{equation}
 \cR_{\ell,m}[H_1^{-1}]
 =\kk\bigl[
 h_1^{\pm1},\ldots,h_\ell^{\pm1},
 \xi_{r;a,b}:2\le r\le m,
 1\le a\le b\le\ell
 \bigr].
 \label{eq:localized-ring-xi}
\end{equation}

\subsection{Polynomial lifts of the Gauss coordinates}
\label{subsec:gaussian-coordinate-invariants}

The functions $\xi_{r;a,b}$ are regular on the quotient of
$X_{\ell,m}^{\mathrm G}$ but are generally rational on $X_{\ell,m}$.  We
define polynomial lifts of these coordinates.  Put $I_a=\{1,\ldots,a\}$ and, for
$1\le p\le a\le\ell$, define
\[ \gamma_{a,p}(C_1) :=(-1)^{a+p}\det C_1[I_{a-1},I_a\setminus\{p\}], \]
where the determinant of the empty matrix is $1$.  For
$2\le r\le m$ and $1\le a\le b\le\ell$, set
\begin{equation}
 \eta_{r;a,b}
 :=\sum_{p=1}^a\sum_{q=1}^b
 \gamma_{a,p}(C_1)\gamma_{b,q}(C_1)(C_r)_{pq}.
 \label{eq:eta-definition}
\end{equation}

\begin{proposition}
\label{prop:eta-gaussian-coordinate}
Each $\eta_{r;a,b}$ is a polynomial element of $\cR_{\ell,m}$.  On the
Gauss chart,
\begin{equation}
 \eta_{r;a,b}=h_{a-1}h_{b-1}\xi_{r;a,b}.
 \label{eq:eta-slice-formula}
\end{equation}
Consequently,
\begin{equation}
 \cR_{\ell,m}[H_1^{-1}]
 =\kk\bigl[
 h_1^{\pm1},\ldots,h_\ell^{\pm1},
 \eta_{r;a,b}:2\le r\le m,
 1\le a\le b\le\ell
 \bigr].
 \label{eq:localized-ring-eta}
\end{equation}
\end{proposition}

\begin{proof}
Let $C_1=u^TDu$ be \eqref{eq:symmetric-gaussian-decomposition}.  Cramer's
rule gives
\begin{equation}
 (u^{-1})_{p,a}=\frac{\gamma_{a,p}(C_1)}{h_{a-1}},
 \qquad 1\le p\le a.
 \label{eq:u-inverse-cofactor}
\end{equation}
Since $u^{-1}$ is upper triangular,
\[ \xi_{r;a,b} =\sum_{p=1}^a\sum_{q=1}^b (u^{-1})_{p,a}(C_r)_{pq}(u^{-1})_{q,b}. \]
Substituting \eqref{eq:u-inverse-cofactor} proves
\eqref{eq:eta-slice-formula}.  The expression
\eqref{eq:eta-definition} is polynomial and is invariant on the dense open
set $X_{\ell,m}^{\mathrm G}$; hence it is invariant on all of
$X_{\ell,m}$.  Equation \eqref{eq:localized-ring-eta} follows from
\eqref{eq:localized-ring-xi}.
\end{proof}


\section{Two-sided matrix invariants and symmetric restriction}
\label{sec:matrix-invariants}

Assume $\ell,m\ge2$.  The hive-quiver construction for complete triple
flags appears in \cite{Fsemi}; the flagged Kronecker construction used here
is developed in \cite{FeiKroneckerII}.  We formulate the latter in terms of
two-sided matrix invariants and then restrict it to symmetric matrices.  The
transfer is recalled in Appendix~\ref{app:matrix-transfer}.

\subsection{Matrix invariant algebras}
\label{subsec:two-matrix-invariant-algebras}

Retain $X_{\ell,m}=\SymMat_\ell^m$ from Section~\ref{sec:plethysm-highest-weight}, and put
\[ Y_{\ell,m}:=\Mat_\ell^m, \qquad X_{\ell,m}\subseteq Y_{\ell,m}. \]
We use two copies $U_\ell^{\mathrm L}$ and $U_\ell^{\mathrm R}$ of the upper
unitriangular group, acting on $Y_{\ell,m}$ by
\begin{equation}
 (u,v)\cdot(A_1,\ldots,A_m)
 =\bigl(u^{-T}A_1v^{-1},\ldots,u^{-T}A_mv^{-1}\bigr).
 \label{eq:two-sided-action}
\end{equation}
Define the \emph{two-sided matrix invariant algebra}
\[ \cT_{\ell,m} :=\kk[Y_{\ell,m}]^{U_\ell^{\mathrm L}\times U_\ell^{\mathrm R}}. \]
The diagonal embedding $u\mapsto(u,u)$ preserves $X_{\ell,m}$ and induces
the congruence action used in \eqref{eq:main-ring-matrices}.  Restriction
therefore gives a graded homomorphism
\[ \rho:\cT_{\ell,m}\longrightarrow\cR_{\ell,m}. \]
Transposition defines an involution
\[ \tau_Y(A_1,\ldots,A_m)=(A_1^T,\ldots,A_m^T) \]
of $Y_{\ell,m}$.  It exchanges the two unipotent factors in
\eqref{eq:two-sided-action}, hence induces an involution $\tau$ of
$\cT_{\ell,m}$.  Since $X_{\ell,m}$ is the fixed locus of $\tau_Y$,
one has $\rho\circ\tau=\rho.$
For $1\le r\le m$, put
\[ \Delta_r:=\det A_r\in\cT_{\ell,m}; \]
we use the same notation for its restriction $\det C_r$.

\begin{proposition}
\label{prop:basic-rings}
The algebras $\cT_{\ell,m}$ and $\cR_{\ell,m}$ are normal factorial domains,
and
\begin{align*}
 \dim\cT_{\ell,m}
 &=m\ell^2-\ell(\ell-1)=(m-1)\ell^2+\ell,\\
 \dim\cR_{\ell,m}
 &=m\binom{\ell+1}{2}-\binom{\ell}{2}
   =\frac{(m-1)\ell^2+(m+1)\ell}{2}.
\end{align*}
\end{proposition}

\begin{proof}
Invariant subrings of normal domains are normal.  These maximal-unipotent
invariant rings are finitely generated by Grosshans' theorem.  Factoriality
follows from Popov's theorem, because the acting unipotent groups are
connected and have no nontrivial algebraic characters; see
\cite[Theorem~3.17]{PopovVinberg} and \cite{Grosshans}.

For $\cT_{\ell,m}$, the point $(I_\ell,0,\ldots,0)$ has trivial stabilizer:
$u^{-T}v^{-1}=I_\ell$ forces the upper unitriangular matrix $v$ to equal the
lower unitriangular matrix $u^{-T}$, so $u=v=I_\ell$.  For $\cR_{\ell,m}$,
the same point has trivial stabilizer under congruence, since
$u^{-T}u^{-1}=I_\ell$ forces $u=I_\ell$.  The formulas now follow from Rosenlicht's dimension theorem
\cite{Rosenlicht}.
\end{proof}

The torus $T_\ell^{\mathrm L}\times T_\ell^{\mathrm R}\times T_W$ grades
$\cT_{\ell,m}$.  If a homogeneous function has degree
$(\lambda^{\mathrm L},\lambda^{\mathrm R};\alpha)$, then its symmetric
restriction has degree
\begin{equation}
 (\lambda^{\mathrm L}+\lambda^{\mathrm R};\alpha).
 \label{eq:degree-restriction}
\end{equation}

\subsection{The two-sided Gauss chart}
\label{subsec:two-sided-gaussian-slice}

The two-sided action has an analogous Gauss chart.  For the first
matrix in $Y_{\ell,m}$, set
\[ \widetilde h_i:=\det A_1[I_i,I_i], \qquad \widetilde H_1:=\prod_{i=1}^{\ell}\widetilde h_i, \]
and let $Y_{\ell,m}^{\mathrm G}$ be the principal open subset
$\widetilde H_1\ne0$.

\begin{proposition}[two-sided Gauss chart]
\label{prop:two-sided-gaussian-slice}
The morphism
\begin{align}
 \widetilde\Psi:
 T_\ell\times\Mat_\ell^{m-1}\times U_\ell\times U_\ell
 &\longrightarrow Y_{\ell,m}^{\mathrm G},
 \notag\\
 (D,B_2,\ldots,B_m,u,v)
 &\longmapsto
 (u^TDv,u^TB_2v,\ldots,u^TB_mv)
 \label{eq:two-sided-gaussian-map}
\end{align}
is an isomorphism.  The action of $(a,b)\in
U_\ell^{\mathrm L}\times U_\ell^{\mathrm R}$ is right translation on the
last two factors:
\begin{equation}
 (a,b)\cdot(D,B_2,\ldots,B_m,u,v)
 =(D,B_2,\ldots,B_m,ua^{-1},vb^{-1}).
 \label{eq:two-sided-gaussian-action}
\end{equation}
Consequently,
\begin{equation}
 \cT_{\ell,m}[\widetilde H_1^{-1}]
 \cong\kk[T_\ell\times\Mat_\ell^{m-1}].
 \label{eq:two-sided-gaussian-quotient}
\end{equation}
Under this identification, transposition acts by
\begin{equation}
 (D,B_2,\ldots,B_m)
 \longmapsto(D,B_2^T,\ldots,B_m^T).
 \label{eq:transpose-on-two-sided-slice}
\end{equation}
Its fixed locus is therefore
$T_\ell\times\SymMat_\ell^{m-1}$, which is the quotient in
Proposition~\ref{prop:gaussian-slice}.
\end{proposition}

\begin{proof}
The nonvanishing of the leading principal minors is equivalent to the
existence of a unique $LDU$ decomposition
\[ A_1=u^TDv, \qquad u,v\in U_\ell, \qquad D\in T_\ell. \]
The remaining normalized matrices are $B_r=u^{-T}A_rv^{-1}$.  Gaussian
elimination shows that these assignments and their inverses are regular on
$Y_{\ell,m}^{\mathrm G}$, proving
\eqref{eq:two-sided-gaussian-map}.  Formula
\eqref{eq:two-sided-gaussian-action} follows from the two-sided action.
Taking invariants removes the last two factors and proves
\eqref{eq:two-sided-gaussian-quotient}.  Finally,
\[ (u^TDv,u^TB_2v,\ldots,u^TB_mv)^T =(v^TDu,v^TB_2^Tu,\ldots,v^TB_m^Tu), \]
which gives \eqref{eq:transpose-on-two-sided-slice}.  The assertion about
the fixed locus follows.
\end{proof}

\begin{corollary}
\label{cor:localized-restriction-surjective}
Restriction to symmetric matrices induces a surjection
\[ \rho_{H_1}:\cT_{\ell,m}[\widetilde H_1^{-1}] \twoheadrightarrow\cR_{\ell,m}[H_1^{-1}]. \]
\end{corollary}

\begin{proof}
Under the two Gauss-chart identifications, this is the surjective
restriction
\[ \kk[T_\ell\times\Mat_\ell^{m-1}] \longrightarrow\kk[T_\ell\times\SymMat_\ell^{m-1}]. \]
\end{proof}


\subsection{Transfer from the flagged Kronecker model}
\label{subsec:matrix-transfer}

The construction in \cite{FeiKroneckerII} is stated for the semi-invariant
ring of a flagged Kronecker representation space.  Appendix~\ref{app:matrix-transfer}
identifies that ring with $\cT_{\ell,m}$ and compares the corresponding
cluster functions.

\begin{proposition}\label{prop:matrix-transfer}
Let $\mathsf K_{\ell,m}$ be the $\ell$-flagged $m$-arrow Kronecker quiver
with its standard dimension vector $\boldsymbol\beta_\ell$.  There is a
natural multigraded algebra isomorphism
\[ \SI_{\boldsymbol\beta_\ell}(\mathsf K_{\ell,m}) \xrightarrow{\sim}\cT_{\ell,m}. \]
It sends the determinant of the $r$-th central arrow to $\Delta_r$ and
intertwines reflection of the two flag arms with matrix transposition.
\end{proposition}

The proof is given in Appendix~\ref{app:matrix-transfer}.

\subsection{Alternating words and chamber determinants}
\label{subsec:matrix-chambers}

For $0\le a\le\ell$, let
\[ J_a:\kk^a\hookrightarrow\kk^\ell, \qquad P_a:=J_a^T:\kk^\ell\twoheadrightarrow\kk^a \]
be the standard inclusion and projection.  Empty blocks are omitted.  If
$I=(r_1,\ldots,r_{2q+1})$ has odd length, define on $(\GL_\ell)^m$
\begin{equation}
 A_I:=A_{r_1}A_{r_2}^{-1}A_{r_3}\cdots
       A_{r_{2q}}^{-1}A_{r_{2q+1}},
 \label{eq:alternating-word}
\end{equation}
and let $\overleftarrow I=(r_{2q+1},\ldots,r_1)$.  Put
\[
 [a]:=(1,2,\ldots,a),
 \qquad
 \mathbf q_n:=
 \begin{cases}
  [n-1],&n\text{ even},\\
  [n],&n\text{ odd},
 \end{cases}
 \qquad 2\le n\le m,
\]
and
\[ \Pi_n:=\prod_{\substack{2\le r<n\\r\text{ even}}}\Delta_r. \]

Let $i,j\ge0$ and $k=i+j$.  On $(\GL_\ell)^m$ set
\begin{equation}
 \phi_{i,j}^{(n),\mathrm L,\circ}(A)
 :=\det
 \begin{bmatrix}
  P_iA_nJ_k\\[1mm]
  P_jA_{\mathbf q_n}J_k
 \end{bmatrix}.
 \label{eq:left-chamber-determinant}
\end{equation}
For $1\le k\le\ell$, define recursively
\begin{equation}
 \phi_{i,j}^{(n),\mathrm L}:=
 \begin{cases}
  \phi_{i,0}^{(n),\mathrm L,\circ},&j=0,\\[1mm]
  \Pi_n\phi_{i,j}^{(n),\mathrm L,\circ},
    &j>0\text{ and }(k<\ell\text{ or }n\text{ even}),\\[1mm]
  \phi_{i,j}^{(n-1),\mathrm L},
    &j>0,\ k=\ell,\ n\text{ odd},
 \end{cases}
 \label{eq:regular-left}
\end{equation}
and set
\begin{equation}
 \phi_{i,j}^{(n),\mathrm R}:=\tau
 \bigl(\phi_{i,j}^{(n),\mathrm L}\bigr).
 \label{eq:regular-right}
\end{equation}
When the middle line of \eqref{eq:regular-left} applies,
\begin{equation}
 \phi_{i,j}^{(n),\mathrm R}(A)
 =\Pi_n\det\left[
  P_kA_nJ_i\ \middle|\ P_kA_{\overleftarrow{\mathbf q_n}}J_j
 \right].
 \label{eq:right-formula}
\end{equation}

\begin{lemma}\label{lem:denominator-clearing}
The functions $\phi_{i,j}^{(n),\mathrm L}$ and
$\phi_{i,j}^{(n),\mathrm R}$ extend to polynomial functions on
$Y_{\ell,m}$.
\end{lemma}

\begin{proof}
Apply Laplace expansion along the top $i$ rows in $\phi_{i,j}^{(n),\mathrm L,\circ}$, expand $\det(A_{\mathbf q_n}[I,J])$ by Cauchy--Binet formula and employ Jacobi's complementary-minor identity
\[ \det(B)\det(B^{-1}[I,J]) =(-1)^{\sum I+\sum J}\det(B[J^c,I^c]). \]
The inverse factor in each term is a product of minors of $A_r^{-1}$ with $r$ even, and $\Pi_n$ supplies their determinants. This shows that $\Pi_n\phi_{i,j}^{(n),\mathrm L,\circ}$ is polynomial.  The
right chamber function is its transpose.
\end{proof}

The alternating words are therefore notation for polynomial functions on
$Y_{\ell,m}$; no determinant is inverted in the coefficient ring.

\begin{lemma}
\label{lem:chamber-invariance}
Every $\phi_{i,j}^{(n),\mathrm L}$ and
$\phi_{i,j}^{(n),\mathrm R}$ belongs to $\cT_{\ell,m}$.
\end{lemma}

\begin{proof}
For every odd-length word $I$,
\[ ((u,v)\cdot A)_I=u^{-T}A_Iv^{-1}. \]
The two block matrices in the chamber determinant are consequently changed
by unitriangular row and column operations.  Their determinants, and the
factors $\Pi_n$, are invariant.  Transposition gives the right chamber
functions.
\end{proof}

\subsection{The flagged Kronecker seed in matrix form}
\label{subsec:flagged-matrix-form}

The determinant formulas above determine the relevant functions up to
nonzero scalars.  The seed of \cite{FeiKroneckerII} fixes a coherent
normalization.

\begin{theorem}\label{thm:flagged-matrix-form}
There exist nonzero scalars $c_{i,j}^{(n)}$ such that, after setting
\[ x_{i,j}^{(n),\mathrm L}:= c_{i,j}^{(n)}\phi_{i,j}^{(n),\mathrm L}, \qquad x_{i,j}^{(n),\mathrm R}:=\tau \bigl(x_{i,j}^{(n),\mathrm L}\bigr), \]
the determinant functions $\Delta_r$ and all distinct chamber functions
$x_{i,j}^{(n),\mathrm L},x_{i,j}^{(n),\mathrm R}$ form the extended cluster
of a full-row-rank seed
\[ \widetilde\Sigma_{\ell,m} =(\widetilde B_{\ell,m},\widetilde{\mathbf x}_{\ell,m}) \quad\text{in }\Frac(\cT_{\ell,m}), \]
and
\begin{equation}
 \cT_{\ell,m}=\mathcal U(\widetilde\Sigma_{\ell,m}),
 \label{eq:flagged-UCA}
\end{equation}
with frozen variables treated as polynomial coefficients.  The normalization
may be chosen so that
\begin{align}
 x_{i,0}^{(n),\mathrm L}&=x_{i,0}^{(n),\mathrm R}
 =\det(P_iA_nJ_i),
 \notag\\
 x_{0,j}^{(2),\mathrm L}&=x_{0,j}^{(2),\mathrm R}
 =\det(P_jA_1J_j),
 \notag\\
 x_{0,j}^{(n),\epsilon}&=x_{0,j}^{(n-1),\epsilon}
 &&(n\ge4\text{ even}),
 \label{eq:even-identification}\\
 x_{i,\ell-i}^{(n),\epsilon}
 &=x_{i,\ell-i}^{(n-1),\epsilon}
 &&(n\ge3\text{ odd}),
 \label{eq:odd-identification}
\end{align}
where $\epsilon\in\{\mathrm L,\mathrm R\}$.  The redundant corners are
normalized as
\[ x_{\ell,0}^{(n),\epsilon}=\Delta_n, \qquad x_{0,\ell}^{(n),\epsilon} =\prod_{\substack{1\le r<n\\r\text{ odd}}}\Delta_r. \]
Transposition preserves the mutable and frozen subsets and satisfies
$\widetilde b_{\tau u,\tau v}=\widetilde b_{uv}.$
\end{theorem}

\begin{proof}
Via Proposition~\ref{prop:matrix-transfer}, the functions above are the
minimally lifted cluster functions of
\cite[Sections~4--5]{FeiKroneckerII}.  The equality
\eqref{eq:flagged-UCA} is \cite[Theorem~0.1]{FeiKroneckerII}, and full row
rank is proved in \cite[Section~5]{FeiKroneckerII}.  Reflection of the two
flags becomes
matrix transposition, including the determinant coefficients introduced by
the minimal lift.
\end{proof}

We use these normalizations from now on.  The notation $f\doteq g$ means
that $f$ and $g$ differ by a nonzero scalar.  Formulas
\eqref{eq:left-chamber-determinant}--\eqref{eq:right-formula} determine the
normalized functions up to $\doteq$; the exchange relations are normalized
with coefficient $+1$.

\subsection{Symmetric restriction and the folded initial family}
\label{subsec:folded-family}

For a symmetric tuple $C$ and any odd-length sequence $I$,
\begin{equation}
 C_I^T=C_{\overleftarrow I}.
 \label{eq:symmetric-word}
\end{equation}
The two determinant matrices in \eqref{eq:left-chamber-determinant} and
\eqref{eq:right-formula} are therefore transposes.

\begin{proposition}\label{prop:transpose-comparison}
For every admissible $n,i,j$, we have
\begin{equation}
 \rho\bigl(x_{i,j}^{(n),\mathrm L}\bigr)
 =\rho\bigl(x_{i,j}^{(n),\mathrm R}\bigr).
 \label{eq:transpose-comparison}
\end{equation}
\end{proposition}

\begin{proof}
By \eqref{eq:symmetric-word}, the two chamber matrices are transposes on the
symmetric locus.
\end{proof}

Denote the common value in \eqref{eq:transpose-comparison} by
\[ z_{i,j}^{(n)} :=\rho\bigl(x_{i,j}^{(n),\mathrm L}\bigr) =\rho\bigl(x_{i,j}^{(n),\mathrm R}\bigr). \]
Let
\[ \mathscr H_\ell :=\{(i,j)\in\mathbb Z_{\ge0}^2:1\le i+j\le\ell\} \setminus\{(\ell,0),(0,\ell)\}. \]
Define
\begin{equation}
 \mathscr V_{\ell,m}
 :=\left(\bigsqcup_{n=2}^m\{n\}\times\mathscr H_\ell\right)\big/\sim,
 \label{eq:V-index}
\end{equation}
where $\sim$ is generated by
\begin{align*}
 (n;0,j)&\sim(n-1;0,j)
 &&(n\ge4\text{ even}),\\
 (n;i,\ell-i)&\sim(n-1;i,\ell-i)
 &&(n\ge3\text{ odd}).
\end{align*}
For $v=[n;i,j]\in\mathscr V_{\ell,m}$, put $z_v=z_{i,j}^{(n)}$, and set
\[ \mathbf x_{\ell,m} :=\{\Delta_1,\ldots,\Delta_m\} \cup\{z_v:v\in\mathscr V_{\ell,m}\}. \]

The principal minors from \eqref{eq:hi-definition} occur in the folded initial
family as
\[ h_i=z_{0,i}^{(2)},\qquad h_\ell=\Delta_1. \]

We write
\[ \omega_a=(\underbrace{1,\ldots,1}_{a},0,\ldots,0) \quad(0\le a\le\ell), \qquad \omega_0=0, \]
for the fundamental polynomial weights of $\GL_\ell$.  We write
$\varepsilon_1,\ldots,\varepsilon_\ell$ for the standard characters of
$T_\ell$, and let $\delta_1,\ldots,\delta_m$ be the standard basis of
$\mathbb Z^m$.  In particular,
\[ \wt_V(\Delta_r)=2\omega_\ell, \qquad \deg_W(\Delta_r)=\ell\delta_r. \]

For $v\in\mathscr V_{\ell,m}$, define its \emph{minimal block index}
\[ s(v):=\min\{n:(n;i,j)\in v\text{ for some }i,j\}. \]
Equivalently, a representative $(n;i,j)$ has minimal block index $n-1$ when
$n\ge4$ is even and $i=0$, or when $n\ge3$ is odd, $i+j=\ell$, and
$j>0$; in every other case its minimal block index is $n$.  Put
$q(v)=\lfloor(s(v)-1)/2\rfloor$ when $j>0$, and $q(v)=0$ when $j=0$.

\begin{proposition}
\label{prop:weights-count}
Let $v=[n;i,j]\in\mathscr V_{\ell,m}$, put $s=s(v)$ and $k=i+j$.  Then
\begin{equation}
 \wt_V(z_v)=\omega_i+\omega_j+\omega_k+2q(v)\omega_\ell.
 \label{eq:V-weight}
\end{equation}
Its $T_W$-multidegree is
\begin{equation}
 \deg_W(z_v)=
 \begin{cases}
  i\delta_s,&j=0,\\[1mm]
  i\delta_s
  +j\displaystyle\sum_{\substack{r<s\\r\text{ odd}}}\delta_r
  +(\ell-j)\displaystyle\sum_{\substack{r<s\\r\text{ even}}}\delta_r,
  &j>0,\ s\text{ even},\\[4mm]
  (i+j)\delta_s
  +j\displaystyle\sum_{\substack{r<s\\r\text{ odd}}}\delta_r
  +(\ell-j)\displaystyle\sum_{\substack{r<s\\r\text{ even}}}\delta_r,
  &j>0,\ s\text{ odd}.
 \end{cases}
 \label{eq:W-degree}
\end{equation}
The cardinality is
\[ |\mathbf x_{\ell,m}| =\frac{(m-1)\ell^2+(m+1)\ell}{2} =\dim\cR_{\ell,m}. \]
\end{proposition}

\begin{proof}
Before restriction, the left determinant uses the first $i$ and $j$ left
rows and the first $k$ right columns.  Each factor in $\Pi_s$
contributes one full left and one full right fundamental weight.  Adding the
left and right weights under \eqref{eq:degree-restriction} gives
\eqref{eq:V-weight}.  Formula \eqref{eq:W-degree} follows by reading the
alternating exponents and then adding the degrees of $\Pi_s$.

Finally,
\[ |\mathscr H_\ell|=\frac{\ell^2+3\ell-4}{2}. \]
There are $m-1$ blocks and $m-2$ boundary identifications, each identifying $\ell-1$ pairs.
After adjoining the $m$ determinant functions, one obtains
\[ m+(m-1)|\mathscr H_\ell|-(m-2)(\ell-1) =\frac{(m-1)\ell^2+(m+1)\ell}{2}. \]
The last equality with the dimension follows from
Proposition~\ref{prop:basic-rings}.
\end{proof}

Theorem~\ref{thm:folded-family-transcendence-basis} shows that
$\mathbf x_{\ell,m}$ is a transcendence basis.  The first two blocks are
given, up to normalization, by
\[ z_{i,j}^{(2)}\doteq \det\left[P_{i+j}C_2J_i\ \middle|\ P_{i+j}C_1J_j\right] \]
and, for $j>0$ and $i+j<\ell$,
\[ z_{i,j}^{(3)}\doteq \Delta_2\det\left[ P_{i+j}C_3J_i\ \middle|\ P_{i+j}(C_3C_2^{-1}C_1)J_j \right]. \]
The second formula is polynomial by Lemma~\ref{lem:denominator-clearing}.

\subsection{Algebraic independence of the folded initial family}
\label{subsec:folded-family-independence}

\begin{theorem}
\label{thm:folded-family-transcendence-basis}
The folded initial family $\mathbf x_{\ell,m}$ is a transcendence basis of
$\Frac(\cR_{\ell,m})$.  In particular,
\[ \Frac(\cR_{\ell,m})=\kk(\mathbf x_{\ell,m}). \]
\end{theorem}

\begin{proof}
Let $\widetilde{\mathbb T}_{\ell,m}$ be the initial cluster torus of
\cite{FeiKroneckerII}.  It is
contained in the Gauss open set because the $n=2$ boundary contains
$\widetilde h_1,\ldots,\widetilde h_\ell$.  Transposition exchanges paired
chamber coordinates and fixes the remaining coordinates.  Its fixed
subscheme is therefore a torus with one coordinate for each transpose
orbit.

By Proposition~\ref{prop:two-sided-gaussian-slice}, the fixed locus of the
two-sided Gauss quotient is
\[ T_\ell\times\SymMat_\ell^{m-1} =\operatorname{Spec}\cR_{\ell,m}[H_1^{-1}]. \]
The fixed torus is a nonempty open subset of this irreducible variety, and
its coordinates restrict to $\mathbf x_{\ell,m}$ by
Proposition~\ref{prop:transpose-comparison}.  Hence
$\Frac(\cR_{\ell,m})=\kk(\mathbf x_{\ell,m}). $
Together with $|\mathbf x_{\ell,m}|=\dim\cR_{\ell,m}$, this proves
algebraic independence.
\end{proof}

\section{Folding the matrix seed}
\label{sec:folding}

We fold the matrix seed of Theorem~\ref{thm:flagged-matrix-form} by
transposition.  Theorem~\ref{thm:folded-family-transcendence-basis} has
already identified the folded cluster functions as a transcendence basis.
The exchange matrix and exchange relations are determined below.  We use
the exchange-matrix conventions of \cite{BFZ}, with frozen variables as
polynomial coefficients, and the standard admissible-orbit terminology of
\cite{DupontFolding}.

\subsection{Folding after restriction}
\label{subsec:restriction-folding}

We use extended exchange matrices with mutable rows and all variables as
columns, with convention
\[ x_ux_u' =\prod_vx_v^{[b_{uv}]_+} +\prod_vx_v^{[-b_{uv}]_+}. \]
Let $\widetilde I$ be the index set of an extended cluster, and let $\tau$ be
an involution preserving its mutable subset.  We say that $\tau$ is
\emph{admissible} for $\widetilde B$ if
\begin{align}
 \widetilde b_{\tau u,\tau v}&=\widetilde b_{uv},
 \label{eq:admissible-equivariance}\\
 \widetilde b_{u,\tau u}&=0
 &&(u\text{ mutable}),
 \label{eq:admissible-orbit}\\
 \widetilde b_{uv}\widetilde b_{u,\tau v}&\ge0
 &&(u\text{ mutable}).
 \label{eq:admissible-sign}
\end{align}
For orbits $\bar u,\bar v$, define
\begin{equation}
 b_{\bar u,\bar v}
 :=\sum_{v'\in\bar v}\widetilde b_{u,v'},
 \qquad u\in\bar u.
 \label{eq:orbit-sum}
\end{equation}

\begin{lemma}\label{lem:restriction-folding}
Let $R,S$ be domains and $\rho_0:R\to S$ a homomorphism.  Suppose that
$\widetilde{\mathbf x}$ satisfies the exchange relations of
$\widetilde B$ in $R$, that $\tau$ is admissible, and that
\[ \rho_0(\widetilde x_v)=\rho_0(\widetilde x_{\tau v})=:z_{\bar v}. \]
Then restriction of the relation at $u$ is
\[ z_{\bar u}z_{\bar u}' =\prod_{\bar v}z_{\bar v}^{[b_{\bar u,\bar v}]_+} +\prod_{\bar v}z_{\bar v}^{[-b_{\bar u,\bar v}]_+}, \]
where $z_{\bar u}'=\rho_0(\widetilde x_u')$.  This value is independent of
the representative $u\in\bar u$.
\end{lemma}

\begin{proof}
Condition \eqref{eq:admissible-sign} implies that all entries in a fixed
column orbit have the same sign.  Hence
\[ \sum_{v'\in\bar v}[\widetilde b_{u,v'}]_+ =\left[\sum_{v'\in\bar v}\widetilde b_{u,v'}\right]_+, \]
and similarly for negative parts.  Apply $\rho_0$ to the unfolded exchange
relation.  Equivariance gives representative independence.
\end{proof}


Under the transfer isomorphism, the initial exchange relations of
$\widetilde\Sigma_{\ell,m}$ are those of
\cite[Corollaries~5.3--5.4]{FeiKroneckerII}; their one-step mutations are
regular by Theorem~\ref{thm:flagged-matrix-form}.  The restrictions of these
relations to symmetric matrices are written in
Theorem~\ref{thm:folded-relations}.

\subsection{Admissibility and the folded matrix}
\label{subsec:folded-matrix}

\begin{proposition}
\label{prop:transpose-admissible}
Transposition is admissible for $\widetilde B_{\ell,m}$.
\end{proposition}

\begin{proof}
Transposition exchanges the left and right interior and boundary-identification relations and
fixes the common-boundary relations.  The determinant coefficients are
fixed and occur in transpose-paired positions, proving
\eqref{eq:admissible-equivariance}.  Distinct members of a nontrivial transpose orbit are never adjacent.  In
an interior or boundary-identification relation at most one member of a nontrivial orbit occurs
in either monomial, while in a common-boundary relation both members occur
in the same monomial.  Hence \eqref{eq:admissible-orbit} and
\eqref{eq:admissible-sign} hold.
\end{proof}

Only initial admissibility is needed for the upper-cluster argument.
In Section~\ref{sec:reddening}, fixed diamond diagonals are removed directly
in the folded matrix.  The unfolded construction is used only for pairs of
transpose-related disk flips with disjoint supports; admissibility is
preserved along these sequences.

Introduce determinant indices $\partial_1,\ldots,\partial_m$ and put
$z_{\partial_r}=\Delta_r$.  Let
\[ \mathscr I_{\ell,m} :=\{\partial_1,\ldots,\partial_m\}\sqcup\mathscr V_{\ell,m}. \]
The outer boundary is
\begin{equation*}
\mathscr B_{\ell,m}:=
 \begin{cases}
  \{[m;0,j]:1\le j\le\ell-1\},&m\text{ odd},\\
  \{[m;i,\ell-i]:1\le i\le\ell-1\},&m\text{ even}.
 \end{cases}
\end{equation*}
Set
\[ \mathscr F_{\ell,m} :=\{\partial_1,\ldots,\partial_m\}\sqcup\mathscr B_{\ell,m}, \qquad \mathscr M_{\ell,m}:=\mathscr I_{\ell,m}\setminus\mathscr F_{\ell,m}. \]
Then
\[ |\mathscr F_{\ell,m}|=m+\ell-1, \qquad |\mathscr M_{\ell,m}| =\frac{(m-1)(\ell-1)(\ell+2)}2. \]

The complete quiver of the first two-block example is shown in
Figure~\ref{fig:seed-43}.

\begin{figure}[t]
\centering
\Symfourthree
\caption{The complete folded seed $\Sigma_{4,3}$.  The upper and lower
triangles are the blocks $n=2$ and $n=3$; the row $31,22,13$ is their
identified boundary.  A label $ij$ in block $n$ denotes
$z_{i,j}^{(n)}$, while a label on the common row denotes
$z_{i,4-i}^{(2)}=z_{i,4-i}^{(3)}$.  Circles are mutable and squares are
frozen.  The ordered arrow valuation is read from source to target; for a
frozen endpoint it is determined by the orbit-size symmetrizer.  The red
and blue valued arrows use the color convention of the two-matrix figure,
and the two parallel black arrows at the top have valuation $(2,2)$.}
\label{fig:seed-43}
\end{figure}

For odd $n\ge3$, block $n$ is identified with block $n-1$ along the side
$i+j=\ell$; for even $n\ge4$, the identification is along the side $i=0$.
A chamber fixed by transposition gives an orbit of size one; all other
chamber labels occur in transpose pairs.  The remaining boundary of block
$m$, together with $\Delta_1,\ldots,\Delta_m$, is frozen.

Define
\[ B_{\ell,m}=(b_{uv})_{ \substack{u\in\mathscr M_{\ell,m}\\v\in\mathscr I_{\ell,m}}} \]
by the orbit-sum formula \eqref{eq:orbit-sum}.

\begin{proposition}
\label{prop:folded-B}
The matrix $B_{\ell,m}$ is well-defined and has full row rank.  Its mutable
principal part is skew-symmetrizable.  If $d_u$ is the size
of the transpose orbit above $u$, then
\begin{equation}
 d_ub_{uv}=-d_vb_{vu}
 \qquad(u,v\in\mathscr M_{\ell,m}).
 \label{eq:symmetrizer}
\end{equation}
\end{proposition}

\begin{proof}
Well-definedness follows from equivariance.  For mutable $u,v$,
\[ d_ub_{uv} =\sum_{u'\in\bar u}\sum_{v'\in\bar v} \widetilde b_{u'v'} =-d_vb_{vu}, \]
which proves \eqref{eq:symmetrizer}.

View $\widetilde B_{\ell,m}$ as a surjective, $\tau$-equivariant linear map
over $\mathbb Q$ from the column space to the mutable row space.  Taking
invariants under a group of order two is exact in characteristic zero.  In
orbit-sum bases, the induced map is represented by $B_{\ell,m}$ up to
invertible diagonal rescaling by orbit sizes.  Hence $B_{\ell,m}$ is
surjective and has full row rank.
\end{proof}

\subsection{The folded exchange relations}
\label{subsec:folded-relations}

We use the conventions
\[ z_{0,0}^{(n)}:=1, \qquad z_{\ell,0}^{(n)}:=\Delta_n, \qquad z_{0,\ell}^{(2)}:=\Delta_1. \]
All variables on identified boundaries are interpreted through the identifications in
\eqref{eq:V-index}.

\begin{theorem}[folded exchange relations]
\label{thm:folded-relations}
For every mutable index, the mutation prescribed by $B_{\ell,m}$ is a
homogeneous regular element of $\cR_{\ell,m}$.  The relations are:

\smallskip
\noindent\emph{(a) Variables on the base edge.}  For $2\le n\le m$ and
$1\le i\le\ell-2$,
\begin{equation}
 z_{i,0}^{(n)}(z_{i,0}^{(n)})'
 =z_{i+1,0}^{(n)}(z_{i-1,1}^{(n)})^2
  +z_{i-1,0}^{(n)}(z_{i,1}^{(n)})^2.
 \label{eq:base-edge}
\end{equation}
At the corner $i=\ell-1$,
\begin{align}
 z_{\ell-1,0}^{(n)}(z_{\ell-1,0}^{(n)})'
 &=\Delta_n(z_{\ell-2,1}^{(n)})^2
   +z_{\ell-2,0}^{(n)}(z_{\ell-1,1}^{(n)})^2,
 &&n\text{ even},
 \label{eq:base-edge-even-corner}\\
 z_{\ell-1,0}^{(n)}(z_{\ell-1,0}^{(n)})'
 &=(z_{\ell-2,1}^{(n)})^2
   +\Delta_nz_{\ell-2,0}^{(n)}(z_{\ell-1,1}^{(n)})^2,
 &&n\text{ odd}.
 \label{eq:base-edge-odd-corner}
\end{align}

\smallskip
\noindent\emph{(b) The other boundary of the block $n=2$.}  For
$1\le j\le\ell-1$,
\begin{equation}
 z_{0,j}^{(2)}(z_{0,j}^{(2)})'
 =z_{0,j-1}^{(2)}(z_{1,j}^{(2)})^2
  +z_{0,j+1}^{(2)}(z_{1,j-1}^{(2)})^2.
 \label{eq:first-boundary}
\end{equation}

\smallskip
\noindent\emph{(c) Chamber interiors.}  For $2\le n\le m$, $i,j>0$, and
$i+j<\ell$,
\begin{align}
 z_{i,j}^{(n)}(z_{i,j}^{(n)})'
 & =z_{i,j-1}^{(n)}z_{i+1,j}^{(n)}z_{i-1,j+1}^{(n)}
+z_{i,j+1}^{(n)}z_{i-1,j}^{(n)}z_{i+1,j-1}^{(n)}.
 \label{eq:interior}
\end{align}

\smallskip
\noindent\emph{(d) Even boundary identifications.}  Let $2\le n<m$ be even and let
$i,j\ge1$ with $i+j=\ell$.  Then
\begin{align}
 z_{i,j}^{(n)}(z_{i,j}^{(n)})'
 &=c_{n;i,j}
   z_{i-1,j}^{(n)}z_{i,j-1}^{(n+1)}
+z_{i,j-1}^{(n)}z_{i-1,j}^{(n+1)},
 \label{eq:even-gluing}
\end{align}
where
\begin{equation}
 c_{n;i,j}:=
 \begin{cases}
  \Delta_n,&(i,j)=(\ell-1,1),\\
  1,&\text{otherwise}.
 \end{cases}
 \label{eq:even-corner-factor}
\end{equation}

\smallskip
\noindent\emph{(e) Odd boundary identifications.}  Let $3\le n<m$ be odd and
$1\le j\le\ell-1$.  Then
\begin{align}
 z_{0,j}^{(n)}(z_{0,j}^{(n)})'
 &=\kappa_{n;j}z_{1,j}^{(n)}z_{1,j-1}^{(n+1)}
+z_{1,j-1}^{(n)}z_{1,j}^{(n+1)},
 \label{eq:odd-gluing}
\end{align}
where
\begin{equation}
 \kappa_{n;j}:=
 \begin{cases}
  \Delta_n,&j=\ell-1,\\
  1,&j<\ell-1.
 \end{cases}
 \label{eq:kappa}
\end{equation}
\end{theorem}

\begin{proof}
Apply Lemma~\ref{lem:restriction-folding} to
$\widetilde\Sigma_{\ell,m}$ and $\rho$.  Proposition
\ref{prop:transpose-comparison} identifies the two members of every
transpose pair.  Products of left and right chamber functions on common
bases become squares, giving \eqref{eq:base-edge}--
\eqref{eq:base-edge-odd-corner} and \eqref{eq:first-boundary}.  Interior relations restrict to
\eqref{eq:interior}; the two alternating edge families give
\eqref{eq:even-gluing} and \eqref{eq:odd-gluing}.

The minimal lift in \cite{FeiKroneckerII} inserts a determinant at three
corner types.
At the corner of the common base edge it gives the parity-dependent relations
\eqref{eq:base-edge-even-corner} and
\eqref{eq:base-edge-odd-corner}.  On the other boundary of the block $n=2$,
the correction is encoded by $z_{0,\ell}^{(2)}=\Delta_1$.  The two boundary-identification
corrections are \eqref{eq:even-corner-factor} and \eqref{eq:kappa}.  The unfolded mutated
variables are regular in $\cT_{\ell,m}$ by
Theorem~\ref{thm:flagged-matrix-form}; their restrictions are therefore regular
in $\cR_{\ell,m}$.  Homogeneity follows by restricting the unfolded grading.
\end{proof}

The five families account for
\begin{align*}
 &(m-1)(\ell-1)+(\ell-1)
 +(m-1)\frac{(\ell-1)(\ell-2)}2\\
 &\hspace{30mm}+(m-2)(\ell-1)
 =\frac{(m-1)(\ell-1)(\ell+2)}2
\end{align*}
mutable indices.  Homogeneity also gives
\[ B_{\ell,m}\operatorname{Wt}(\mathbf x_{\ell,m})^T=0, \]
where $\operatorname{Wt}(\mathbf x_{\ell,m})$ is the weight matrix of the
folded family.

\begin{corollary}[the folded seed]
\label{cor:folded-seed}
The pair
$\Sigma_{\ell,m}:=(B_{\ell,m},\mathbf x_{\ell,m})$
is a full-row-rank seed with skew-symmetrizable principal part.  Its
one-step mutations are regular homogeneous invariants.
\end{corollary}

\begin{proof}
Algebraic independence is Theorem~\ref{thm:folded-family-transcendence-basis};
Theorem~\ref{thm:folded-relations} gives the exchange relations, and
Proposition~\ref{prop:folded-B} gives full row rank and
skew-symmetrizability.
\end{proof}

The localization at $H_1$ inverts $\Delta_1=h_\ell$ but no
other central determinant.  The proof of the main theorem will use the
analogous localization for $C_2$ and then intersect the two localized
rings.

\section{The upper cluster algebra}
\label{sec:upper-cluster-equality}

Let
\[ \cU_{\ell,m}:=\mathcal U(\Sigma_{\ell,m}) \]
be the upper cluster algebra of the folded seed, with frozen variables as
polynomial coefficients.  We prove $\cU_{\ell,m}\subseteq\cR_{\ell,m}$ from the initial seed and
its adjacent seeds, then obtain the reverse inclusion
after Gauss localization and remove the localization by intersection.

\subsection{The initial upper bound}
\label{subsec:upper-bound-regularity}

Let
$\mathscr I=\mathscr M_{\ell,m}\sqcup\mathscr F_{\ell,m}$
be the index set of the folded seed.  Its initial Laurent ring is
\[ \mathcal L_0 :=\kk[z_f:f\in\mathscr F_{\ell,m}] [z_u^{\pm1}:u\in\mathscr M_{\ell,m}]. \]
For a mutable index $u$, let $z_u'$ be its one-step mutation and put
\[ \mathcal L_u :=\kk[z_f:f\in\mathscr F_{\ell,m}] [(z_u')^{\pm1},z_v^{\pm1}:v\in\mathscr M_{\ell,m}\setminus\{u\}]. \]
The corresponding upper bound is
\[ \overline{\cU}_{\ell,m}:=\mathcal L_0\cap\bigcap_{u\in\mathscr M_{\ell,m}}\mathcal L_u\subseteq\Frac(\cR_{\ell,m}). \]

By Proposition~\ref{prop:folded-B}, the extended exchange matrix has full
row rank.  The ordinary-cluster specialization of
\cite[Theorem~3.11]{GSV}, over the polynomial ring in the frozen variables,
therefore gives
\begin{equation}
 \cU_{\ell,m}=\overline{\cU}_{\ell,m}.
 \label{eq:upper-bound-equality}
\end{equation}

\subsection{Initial irreducibility}
\label{subsec:initial-irreducibility}

For a polynomial dominant weight
$\lambda=(\lambda_1,\ldots,\lambda_\ell)$, we call $\lambda_1$ its first row.  We use
the classical decomposition
\begin{equation}
 \Sym^d(\Sym^2V)
 \cong\bigoplus_{\rho\vdash d}\mathbf S_{2\rho}V,
 \label{eq:sym-sym2-decomposition}
\end{equation}
with summands of height greater than $\ell$ omitted; see
\cite[Chapter~I, Section~5, Example~4]{Macdonald}.

\begin{lemma}\label{lem:first-row-bound}
Let $0\ne f\in\cR_{\ell,m}$ be homogeneous for $T_V\times T_W$.
Its irreducible factors are homogeneous.  If $f$ is nonconstant and has
$T_V$-weight $\lambda$, then $\lambda_1\ge2$.  Consequently, $f$ is irreducible
whenever $\lambda_1\le3$.
\end{lemma}

\begin{proof}
Since $\cR_{\ell,m}$ is a UFD with constant units, the connected torus
$T_V\times T_W$ fixes each irreducible factor up to scalar; the factors are
therefore homogeneous.  The possible $T_V$-weights in multidegree $\alpha$
are the highest weights occurring in
$\bigotimes_{r=1}^m\Sym^{\alpha_r}(\Sym^2V). $
By \eqref{eq:sym-sym2-decomposition}, every nontrivial factor contributes a
partition whose first row has length at least two, and the same is true for
every Littlewood--Richardson constituent.  Thus a nonconstant homogeneous
element has first row at least two.  A product of two such elements has
first row at least four.
\end{proof}

\begin{corollary}
\label{cor:initial-irreducibility}
The initial functions below are irreducible in $\cR_{\ell,m}$:
\begin{enumerate}
 \item every central determinant $\Delta_r$;
 \item every base-edge function $z_{i,0}^{(n)}$;
 \item every chamber function whose minimal block index is $2$.
\end{enumerate}
\end{corollary}

\begin{proof}
The corresponding weights are
$2\omega_\ell$, $2\omega_i$, and $\omega_i+\omega_j+\omega_{i+j},$
whose first rows have length at most three.
\end{proof}

The remaining cases reduce to determinants of symmetric pairings.

\begin{lemma}\label{lem:symmetric-pairing-determinant}
Let $E$ be a vector space over a field of characteristic different from
two, and let $A,B\subseteq E$ be subspaces of the same positive
dimension $r$.  For $C\in\Sym^2(E^\vee)$, let
\[ p_{A,B}(C):=\det\bigl(C(b_i,a_j)\bigr)_{1\le i,j\le r}, \]
where bases of $A$ and $B$ have been chosen.  Then
$p_{A,B}$ is absolutely irreducible, up to a nonzero scalar depending
on the bases.
\end{lemma}

\begin{proof}
Let $N=\dim\Sym^2(E^\vee)$ and consider
\[ \mathcal I_{A,B} =\{([a],C)\in\mathbb P(A)\times\Sym^2(E^\vee):C(a,B)=0\}. \]
For $a\ne0$, the map $C\mapsto C(a,-)|_B$ from $\Sym^2(E^\vee)$ to
$B^\vee$ is surjective.  Thus $\mathcal I_{A,B}$ is a vector bundle over
$\mathbb P(A)$ of dimension $N-1$, hence irreducible.  Its proper image in
$\Sym^2(E^\vee)$ is the hypersurface $V(p_{A,B})$.

The polynomial $p_{A,B}$ is nonzero: choose bases adapted to
$A\cap B$ and a symmetric form whose pairing between $A$ and $B$ is
nonsingular.  The same bases admit a one-parameter family with pairing
matrix $\operatorname{diag}(1,\ldots,1,t)$.  Hence $V(p_{A,B})$ contains a
smooth corank-one point and is reduced.  It is therefore an irreducible
hypersurface.  The argument is unchanged after field extension, so
$p_{A,B}$ is absolutely irreducible.
\end{proof}

The coefficient argument uses two rank estimates.  For a block index $s\ge3$, work on the open set where
$C_1,\ldots,C_{s-1}$ are invertible.  Write
\begin{equation*}
G_s:=
 \begin{cases}
  C_{\overleftarrow{\mathbf q_s}},&s\text{ even},\\
  C_{s-1}^{-1}C_{s-2}\cdots C_2^{-1}C_1,
    &s\text{ odd},
 \end{cases}
\end{equation*}
where the products alternate as in \eqref{eq:alternating-word}; for odd
$s$, one has $C_{\overleftarrow{\mathbf q_s}}=C_sG_s$.

\begin{lemma}\label{lem:rank-drop-codimension}
Let $k=i+j$.
\begin{enumerate}
 \item If $s$ is even and $i,j>0$, then
 $\operatorname{codim} \{\operatorname{rank}(P_kG_sJ_j)<j\}\ge i+1.$
 \item If $s$ is odd, $j>0$, and $k<\ell$, then
 $\operatorname{codim} \{\operatorname{rank}[J_i\mid G_sJ_j]<k\} \ge\ell-k+1.$
\end{enumerate}
The codimensions are taken in the determinant-open space of the preceding
symmetric matrices.
\end{lemma}

\begin{proof}
Fix the preceding matrices other than $C_1$.  On the determinant-open set
write $G_s=HC_1$ with $H$ invertible.  For every nonzero $x\in\kk^\ell$,
the evaluation map $C\mapsto Cx$ from $\SymMat_\ell$ to $\kk^\ell$ is
surjective.  In the even case, for each $[a]\in\mathbb P^{j-1}$ the
condition
$P_kHC_1J_ja=0$
imposes $k$ independent linear equations.  The incidence variety therefore
has codimension at least $k-(j-1)=i+1$.  In the odd case the condition
$HC_1J_ja\in J_i\kk^i$ imposes $\ell-i$ equations, giving codimension
$(\ell-i)-(j-1)=\ell-k+1$.
\end{proof}

\begin{lemma}
\label{lem:initial-functions-irreducible}
Every member of the folded initial family $\mathbf x_{\ell,m}$ is
irreducible in $\cR_{\ell,m}$.
\end{lemma}

\begin{proof}
Only chamber functions with minimal block index $s\ge3$ remain.  Such a function has
$j>0$; if $s$ is even then $i>0$, while if $s$ is odd then $i+j<\ell$.
Let $A_{<s}$ be the polynomial ring in the entries of
$C_1,\ldots,C_{s-1}$ and $K_{<s}=\Frac(A_{<s})$.

Over $K_{<s}$, the factor $\Pi_s$ is a unit.  If $s$ is even, the block
$P_kG_sJ_j$ has generic rank $j$, and quotienting its column span converts
the chamber determinant into the determinant of $C_s$ on two
$i$-dimensional subspaces.  If $s$ is odd, the matrix
$[J_i\mid G_sJ_j]$ has generic rank $k=i+j$, and the chamber function is the
determinant of $C_s$ on two $k$-dimensional subspaces.  In either case,
Lemma~\ref{lem:symmetric-pairing-determinant} gives irreducibility in
$K_{<s}[\SymMat_\ell]$.

An irreducible common factor of the coefficients would define a divisor in the
rank-drop locus on $\prod_{r<s}\Delta_r\ne0$, contradicting
Lemma~\ref{lem:rank-drop-codimension}.  It must therefore be associated to
some $\Delta_r$, $r<s$.  But \eqref{eq:W-degree} gives degree $j$ or
$\ell-j$ in $C_r$, both smaller than $\deg_{C_r}\Delta_r=\ell$.  Gauss's
lemma proves irreducibility after adjoining $C_s$; adjoining the later
matrix entries preserves primality.  Hence the chamber function is
irreducible in $\kk[X_{\ell,m}]$ and therefore in $\cR_{\ell,m}$.
\end{proof}

\subsection{Coprimality with one-step mutations}
\label{subsec:first-mutation-coprimality}

For a homogeneous element $f$, write
$\deg(f)=(\deg_W(f),\wt_V(f))$.  If
$z_uz_u'=M_{u,+}+M_{u,-} $
is one of the exchange relations of Theorem~\ref{thm:folded-relations},
put
\[ \Gamma_u:=\deg(M_{u,+})-2\deg(z_u). \]
Homogeneity makes the definition independent of the chosen exchange monomial.

\begin{lemma}\label{lem:quotient-degree-obstruction}
If $z_u\mid z_u'$ in $\cR_{\ell,m}$, then the homogeneous component of
$\cR_{\ell,m}$ of bidegree $\Gamma_u$ is nonzero.  In particular, the
$T_W$-part of $\Gamma_u$ is nonnegative and its $T_V$-part is a
polynomial dominant weight.
\end{lemma}

\begin{proof}
The quotient $z_u'/z_u$ would be a nonzero homogeneous element of degree
$\Gamma_u$.
\end{proof}

Substitution of \eqref{eq:V-weight}--\eqref{eq:W-degree} gives
\begin{equation*}
\begin{array}{c|c|c}
  \text{variable}&\text{remaining case}&\text{obstruction otherwise}\\
  \hline
  z_{i,0}^{(n)}&i=1&[\omega_i]\Gamma=-2\ (i\ge2)\\
  z_{0,j}^{(2)}&j=1&[\omega_j]\Gamma=-2\ (j\ge2)\\
  z_{i,j}^{(n)}\text{ interior}&(i,j)=(1,2),(2,1)
    &\text{$T_V$-weight not dominant}\\
  \text{variable on an identified boundary}&\text{none}&[\delta_1]\Gamma=-1
 \end{array}
\end{equation*}
where $[\omega_a]\Gamma$ and $[\delta_r]\Gamma$ denote coefficients in
the indicated bases.

\begin{lemma}\label{lem:degree-reduction}
The quotient-degree obstruction proves $z_u\nmid z_u'$ in every mutable
direction except possibly:
\begin{equation}
 z_{1,0}^{(n)},\qquad z_{0,1}^{(2)},\qquad
 z_{1,2}^{(n)},\qquad z_{2,1}^{(2)}.
 \label{eq:degree-exception-list}
\end{equation}
Here the last two families occur only when their indices are defined.
The two $n=2$ interior exceptions
$z_{1,2}^{(2)}$ and $z_{2,1}^{(2)}$ are also coprime to their mutations.
\end{lemma}

\begin{proof}
For a row on the base edge with $i\ge2$, the coefficient of $\omega_i$ in
the $T_V$-part of $\Gamma_u$ is $-2$.  The same argument on the other
$n=2$ boundary gives coefficient $-2$ at $\omega_j$ for
$j\ge2$.  Thus only $i=1$ and $j=1$, respectively, survive.

For an interior row, put
\[ F(a):=\omega_{a-1}+\omega_{a+1}-\omega_a, \qquad \omega_0:=0. \]
A direct substitution in the weight formula gives
\[ (\Gamma_u)_V=F(i)+F(j)+F(i+j)+c_u\omega_\ell, \]
where $c_u\in\mathbb Z$ records whether a neighboring edge variable is
inherited across an identified boundary.  Since an interior mutable index satisfies
$i+j<\ell$, this correction changes only the coefficient of
$\omega_\ell$ and cannot alter any of the negative coefficients below
${\omega_\ell}$.  The lower coefficients are nonnegative only for $\,(i,j)=(1,2)$ or $(2,1)$.  Assume $i\le j$.  If $i\ge3$, the coefficient at $\omega_i$ is negative
unless $j=i+1$, and in that remaining case the coefficient at
$\omega_{i+j}$ is negative.  If $i=2$, the coefficient at
$\omega_{i+j}$ is negative.  If $i=1$, the coefficient at
$\omega_1$ is nonnegative only when $j=2$.  Interchanging $i$ and
$j$ gives the stated pair.  A direct substitution of the
$T_W$-degrees shows that the $(2,1)$-case has a negative coordinate for
every $n>2$.  In every even or odd boundary-identification row, the coefficient of
$\delta_1$ in $\Gamma_u$ is $-1$.  This proves the reduction to
\eqref{eq:degree-exception-list}.

For $(i,j)=(1,2)$ or $(2,1)$ in the block $n=2$, the putative quotient has
$T_V$-weight
$\omega_2+\omega_4=(2,2,1,1)$
and $T_W$-multidegree $2\delta_1+\delta_2$ or
$\delta_1+2\delta_2$.  Its multiplicity is therefore that of
$\mathbf S_{(2,2,1,1)}V$ in
$\Sym^2(\Sym^2V)\otimes\Sym^2V.$
Since
\[ \Sym^2(\Sym^2V)=\mathbf S_{(4)}V\oplus\mathbf S_{(2,2)}V, \]
the Pieri rule \cite[Chapter~I]{Macdonald} shows that this multiplicity is zero: the partition
$(2,2,1,1)$ is not obtained from either $(4)$ or $(2,2)$ by adding a
horizontal two-strip.  Lemma~\ref{lem:quotient-degree-obstruction} settles
these two cases.
\end{proof}

The degree argument leaves only the three families
\begin{equation}
 z_{1,0}^{(n)},\qquad z_{0,1}^{(2)},\qquad z_{1,2}^{(n)}\quad(n\ge3),
 \label{eq:exceptional-coprimality-families}
\end{equation}
whenever the indicated indices are mutable.

\begin{lemma}\label{lem:transverse-curve-criterion}
Suppose that $z_uz_u'=N_u$ and that a curve
$\gamma:\mathbb A^1\to X_{\ell,m}$ satisfies
\[ \operatorname{ord}_0(z_u\circ\gamma) =\operatorname{ord}_0(N_u\circ\gamma)=1. \]
Then $z_u\nmid z_u'$ in $\cR_{\ell,m}$.
\end{lemma}

\begin{proof}
Write $z_u\circ\gamma=t\eta(t)$ and
$N_u\circ\gamma=tq(t)$ with $\eta(0)q(0)\ne0$.  Then
$z_u'\circ\gamma=q/\eta$ is nonzero at the origin, whereas
$z_u\circ\gamma$ vanishes there.
\end{proof}

\begin{lemma}\label{lem:exceptional-directions}
For every mutable variable in \eqref{eq:exceptional-coprimality-families}, we have
\[ z_u\nmid z_u'. \]
\end{lemma}

\begin{proof}
Appendix~\ref{app:exceptional-coprimality} gives a transverse curve for
each family.  Apply Lemma~\ref{lem:transverse-curve-criterion}.
\end{proof}

\begin{lemma}\label{lem:initial-coprimality}
The initial family consists of pairwise nonassociate prime elements of
$\cR_{\ell,m}$, and for every mutable index $u$,
\begin{equation}
 \gcd(z_u,z_u')=1.
 \label{eq:initial-coprime}
\end{equation}
Consequently,
\begin{equation}
 \cU_{\ell,m}\subseteq\cR_{\ell,m}.
 \label{eq:upper-bound-inclusion}
\end{equation}
\end{lemma}

\begin{proof}
Lemma~\ref{lem:initial-functions-irreducible} and factoriality make every
initial function prime, and algebraic independence makes distinct initial
functions nonassociate.  Lemmas~\ref{lem:degree-reduction} and
\ref{lem:exceptional-directions} prove
\eqref{eq:initial-coprime}.

Let $\nu_u$ be the valuation along $V(z_u)$.  All generators inverted in
$\mathcal L_u$ have $\nu_u$-value zero, so every element of
$\mathcal L_u$ has nonnegative $\nu_u$-valuation.  An element of the upper
bound can have poles in $\mathcal L_0$ only along mutable initial primes,
and membership in the corresponding adjacent Laurent rings removes all
such poles.  Frozen variables are not inverted.  Thus every height-one
valuation of $\cR_{\ell,m}$ is nonnegative on the upper bound. Since a UFD is the intersection
of its height-one valuation rings, the upper bound lies in
$\cR_{\ell,m}$.  Together with \eqref{eq:upper-bound-equality}, this proves the claim.
\end{proof}

\subsection{Restriction and Gauss localizations}
\label{subsec:restriction-upper-bound}

For a mutable folded index $u$, choose a representative in the unfolded seed
$\widetilde u$.  If its transpose orbit has two members, define the
\emph{orbit mutation}
\[ \mu_{\bar u}:=\mu_{\widetilde u}\mu_{\tau\widetilde u}; \]
if the orbit is fixed, use the single mutation $\mu_{\widetilde u}$.  The
two mutations in a nontrivial orbit commute because the two orbit members
are not joined in the unfolded matrix.

\begin{proposition}
\label{prop:source-restriction-in-U}
Restriction to symmetric matrices satisfies
\[ \rho(\cT_{\ell,m})\subseteq\cU_{\ell,m}. \]
\end{proposition}

\begin{proof}
Every $F\in\cT_{\ell,m}=\mathcal U(\widetilde\Sigma_{\ell,m})$ is Laurent
in the unfolded initial seed and in each orbit-mutated seed.  Symmetric
restriction sends the initial Laurent ring to $\mathcal L_0$, and
Lemma~\ref{lem:restriction-folding} sends the orbit-mutated Laurent ring in
direction $u$ to $\mathcal L_u$.  Therefore
\[ \rho(F)\in\mathcal L_0\cap \bigcap_{u\in\mathscr M_{\ell,m}}\mathcal L_u =\cU_{\ell,m}. \]
\end{proof}

For $1\le s\le m$ and $0\le i\le\ell$, set
\[ h_i^{[s]}:=\det C_s[I_i,I_i], \qquad h_0^{[s]}:=1, \qquad H_s:=\prod_{i=1}^{\ell}h_i^{[s]}. \]
Thus $h_i^{[1]}=h_i$ and $h_\ell^{[s]}=\Delta_s$.  Define
$\widetilde h_i^{[s]}$ and $\widetilde H_s$ on $Y_{\ell,m}$ by the
same formulas.

\begin{proposition}
\label{prop:all-localized-restrictions}
For every $s$, restriction induces a surjection
\[ \rho_{H_s}:\cT_{\ell,m}[\widetilde H_s^{-1}] \twoheadrightarrow\cR_{\ell,m}[H_s^{-1}]. \]
\end{proposition}

\begin{proof}
Relabel $C_s$ as the first matrix and apply the two Gauss-chart
propositions.  On the quotient slices the map is the surjective restriction
\[ \kk[T_\ell\times\Mat_\ell^{m-1}] \longrightarrow \kk[T_\ell\times\SymMat_\ell^{m-1}]. \]
\end{proof}

\begin{proposition}
\label{prop:localized-equality}
For every $1\le s\le m$, we have
\[ \cR_{\ell,m}[H_s^{-1}] =\cU_{\ell,m}[H_s^{-1}]. \]
\end{proposition}

\begin{proof}
The functions $h_i^{[s]}$ belong to the initial seed, so $H_s\in\cU_{\ell,m}$.
The inclusion \eqref{eq:upper-bound-inclusion} gives one containment after
localization.  For the other, Proposition~\ref{prop:source-restriction-in-U}
gives
\[ \rho\bigl(\cT_{\ell,m}[\widetilde H_s^{-1}]\bigr) \subseteq\cU_{\ell,m}[H_s^{-1}], \]
and Proposition~\ref{prop:all-localized-restrictions} identifies the
left-hand side with $\cR_{\ell,m}[H_s^{-1}]$.
\end{proof}

\subsection{Removing the Gauss localization}
\label{subsec:remove-gaussian-localization}

We use the first two matrices.  In the first folded chamber,
\[ h_i^{[1]}=z_{0,i}^{(2)}, \qquad h_i^{[2]}=z_{i,0}^{(2)}. \]
For $1\le k\le\ell-1$, denote their exchange polynomials by
\begin{align*}
 E_k^{[1]}
 &:=h_{k-1}^{[1]}(z_{1,k}^{(2)})^2
   +h_{k+1}^{[1]}(z_{1,k-1}^{(2)})^2,\\
 E_k^{[2]}
 &:=h_{k+1}^{[2]}(z_{k-1,1}^{(2)})^2
   +h_{k-1}^{[2]}(z_{k,1}^{(2)})^2.
\end{align*}
Thus
\[ h_k^{[1]}(h_k^{[1]})'=E_k^{[1]}, \qquad h_k^{[2]}(h_k^{[2]})'=E_k^{[2]}. \]

\begin{lemma}
\label{lem:R-localization-intersection}
Inside $\Frac(\cR_{\ell,m})$, we have
\[ \cR_{\ell,m}[H_1^{-1}] \cap\cR_{\ell,m}[H_2^{-1}] =\cR_{\ell,m}. \]
\end{lemma}

\begin{proof}
The factors of $H_1$ and $H_2$ are irreducible by
Lemma~\ref{lem:initial-functions-irreducible}; they are pairwise
nonassociate because their $T_W$-multidegrees are supported on
$\delta_1$ and $\delta_2$, respectively.  Thus $H_1$ and $H_2$
are coprime in the UFD $\cR_{\ell,m}$, and the assertion is the standard
coprime-localization identity.
\end{proof}

\begin{lemma}
\label{lem:U-localization-intersection}
One has
\[ \cU_{\ell,m}[H_1^{-1}] \cap\cU_{\ell,m}[H_2^{-1}] =\cU_{\ell,m}. \]
\end{lemma}

\begin{proof}
Equation~\eqref{eq:upper-bound-equality} expresses $\cU_{\ell,m}$ as a
finite intersection of the initial and adjacent Laurent rings $\mathcal L_v$
(with $v=0$ for the initial ring).  Since a finite intersection is the kernel
of the
difference map from the direct sum to the ambient field, flatness of
localization gives
\begin{equation}
 \cU_{\ell,m}[H_s^{-1}]
 =\bigcap_v\mathcal L_v[H_s^{-1}]
 \qquad(s=1,2).
 \label{eq:localized-upper-bound-star}
\end{equation}
It is enough to prove that $H_1$ and $H_2$ are coprime in every
$\mathcal L_v$.

In the initial ring, and after mutation away from the two boundary
families, $H_1$ and $H_2$ are associated to the polynomial variables
$\Delta_1$ and $\Delta_2$.  After mutation at $h_k^{[1]}$,
\[ H_1\sim\Delta_1E_k^{[1]},\qquad H_2\sim\Delta_2. \]
The polynomial $E_k^{[1]}$ is a sum of two distinct Laurent monomials and
is not divisible by $\Delta_2$.  The analogous mutation at $h_k^{[2]}$
gives
\[ H_1\sim\Delta_1,\qquad H_2\sim\Delta_2E_k^{[2]}, \]
with $E_k^{[2]}$ not divisible by $\Delta_1$.  Each $\mathcal L_v$ is a
UFD, so the coprime-localization identity applies in every Laurent ring in this
intersection.
Intersecting by \eqref{eq:localized-upper-bound-star} proves the claim.
\end{proof}

\begin{theorem}[main theorem]
\label{thm:main-equality}
For all $\ell,m\ge2$, the folded seed $\Sigma_{\ell,m}$ defines the
highest-weight invariant ring as an upper cluster algebra:
\[ \cR_{\ell,m}=\mathcal U(\Sigma_{\ell,m}). \]
\end{theorem}

\begin{proof}
By Proposition~\ref{prop:localized-equality}, the localizations at $H_1$ and
$H_2$ agree.  Lemmas~\ref{lem:R-localization-intersection} and
\ref{lem:U-localization-intersection} therefore give
\[ \cR_{\ell,m} =\cR_{\ell,m}[H_1^{-1}]\cap\cR_{\ell,m}[H_2^{-1}] =\cU_{\ell,m}[H_1^{-1}]\cap\cU_{\ell,m}[H_2^{-1}] =\cU_{\ell,m}. \]
\end{proof}

%Only $\Delta_1$ and $\Delta_2$ are inverted, as factors of $H_1$
%and $H_2$, and both localizations are removed in the proof.






\section{Reddening sequences from a folded reduction}
\label{sec:reddening}

Let
$B_{\ell,m}^{\mathrm{uf}} :=(b_{uv})_{u,v\in\mathscr M_{\ell,m}} $
be the mutable principal part of the folded exchange matrix.  A reddening
sequence for $\Sigma_{\ell,m}$ is a green-to-red sequence for the
principal-coefficient framing of $B_{\ell,m}^{\mathrm{uf}}$; frozen coefficients do not enter the argument.

We use the class $\mathscr P'$ of skew-symmetrizable matrices generated from
the $1\times1$ zero matrix by mutation and source-sink extension.  Equivalently,
a matrix lies in $\mathscr P'$ if and only if it can be reduced to a singleton
by mutations and deletions of sources or sinks.  String-diagram exchange
matrices lie in $\mathscr P'$, and every matrix in $\mathscr P'$ admits a
reddening sequence \cite[Theorem~4.8 and Corollary~4.9]{CaoString}.

To reduce the folded matrix, we remove each transpose-fixed diamond
diagonal directly.  Between two such removals, the required disk flips
occur in disjoint transpose-exchanged regions and descend by orbit mutation.

\subsection{Orbit mutation on separated supports}

Let $\widetilde B$ be a mutable principal matrix carrying an admissible
involution $\tau$, in the sense of
\eqref{eq:admissible-equivariance}--\eqref{eq:admissible-sign}, and let
$B=\operatorname{Fold}_{\tau}(\widetilde B)$ be its orbit-sum fold.

\begin{lemma}\label{lem:orbit-mutation-folding}
For a mutable orbit $\bar k$, the mutations at the members of $\bar k$
commute and
\[ \operatorname{Fold}_{\tau} \left(\prod_{k'\in\bar k}\mu_{k'}(\widetilde B)\right) =\mu_{\bar k}(B). \]
If the members of $\bar k$ are all sources, or all sinks, then $\bar k$ is
a source, respectively a sink, of $B$, and deletion commutes with folding.
\end{lemma}

\begin{proof}
The members of the orbit are not adjacent, so their mutations commute.  For
$\bar i,\bar j\ne\bar k$, sum the simultaneous mutation formula over the
column orbit $\bar j$ and the mutation orbit $\bar k$.  Condition~\eqref{eq:admissible-sign} allows positive and negative
parts to pass through both sums and gives
\[ b'_{\bar i\bar j} =b_{\bar i\bar j} +[b_{\bar i\bar k}]_+[b_{\bar k\bar j}]_+ -[-b_{\bar i\bar k}]_+[-b_{\bar k\bar j}]_+. \]
Entries incident with $\bar k$ change sign on both sides.  The source-sink
statement follows by summing entries of one common sign.
\end{proof}

\begin{lemma}\label{lem:separated-pair}
Suppose the remaining vertices of a $\tau$-equivariant matrix decompose as
\[ \widetilde I=L\sqcup F\sqcup R, \qquad \tau(L)=R, \qquad \tau|_F=\mathrm{id}, \qquad \widetilde B_{L,R}=0. \]
If $\mathbf w=(k_1,\ldots,k_s)$ is supported in $L$, then the paired sequence
$(k_1,\tau k_1),\ldots,(k_s,\tau k_s)$
remains admissible and folds to $(\bar k_1,\ldots,\bar k_s)$.
\end{lemma}

\begin{proof}
A mutation in $L$ cannot create an $L$--$R$ entry, because its mutated
vertex has no neighbor in $R$; the transpose statement holds in $R$.
Thus the zero block and equivariance survive each pair.  A row in $L$ or
$R$ meets at most the orbit member on its own side, while a fixed row meets
the two members equally.  The admissibility signs therefore persist, and
Lemma~\ref{lem:orbit-mutation-folding} applies at every step.
\end{proof}

\subsection{Removing a folded half-diamond diagonal}

Put $N=\ell-1$ and write
\[ d_i^{(r)}=[r;i,0] \qquad(1\le i\le N) \]
for the transpose-fixed diagonal vertices in block $r$.  For even $r$
set
\begin{equation}
 \mathbf p_i^+(r)
 :=\bigl([r;i+u,q-u]\bigr)_{
 {\substack{q=N-i,N-i-1,\ldots,1\\u=0,1,\ldots,q}}},
 \label{eq:positive-removal-sequence}
\end{equation}
and for odd $r$ set
\[ \mathbf p_i^-(r) :=\bigl([r;i-q,j]\bigr)_{ {\substack{q=i-1,i-2,\ldots,1\\j=q,q-1,\ldots,0}}}. \]
The indices are read in the stated order; the corresponding mutations are
performed one anti-diagonal at a time.  Figure~\ref{fig:half-diamond-mutation-order}
shows the order, and Appendix~\ref{app:triangular-mutation} proves the row
formula used below.

\begin{lemma}[removal of a folded half-diamond diagonal]
\label{lem:folded-half-diamond-removal}
Assume that the diagonal labelled $r$ is exposed in the disk model of \cite{FeiKroneckerII}.

If $r$ is even, then, successively for $i=1,\ldots,N$, mutation by
$\mathbf p_i^+(r)$ makes $d_i^{(r)}$ a source or sink, with
\begin{equation}
 \operatorname{supp}b_{d_i^{(r)},\bullet}=
 \begin{cases}
  \{d_{i+1}^{(r)}\},&i<N,\\
  \{[r;N-1,1]\},&i=N,
 \end{cases}
 \qquad
 \begin{cases}
  |b_{d_i^{(r)},d_{i+1}^{(r)}}|=1,&i<N,\\
  |b_{d_N^{(r)},[r;N-1,1]}|=2.&
 \end{cases}
 \label{eq:positive-removal-row}
\end{equation}
Delete $d_i^{(r)}$ and apply the mutations in reverse order.  The matrix on
the surviving vertices is then the preceding matrix with $d_i^{(r)}$ removed.

If $r$ is odd, the same statement holds for $i=N,N-1,\ldots,1$ with
$\mathbf p_i^-(r)$ and
\begin{equation}
 \operatorname{supp}b_{d_i^{(r)},\bullet}=
 \begin{cases}
  \{d_{i-1}^{(r)}\},&i>1,\\
  \{[r;1,1]\},&i=1,
 \end{cases}
 \qquad
 \begin{cases}
  |b_{d_i^{(r)},d_{i-1}^{(r)}}|=1,&i>1,\\
  |b_{d_1^{(r)},[r;1,1]}|=2.&
 \end{cases}
 \label{eq:negative-removal-row}
\end{equation}
\end{lemma}

\begin{proof}
For even $r$, after deleting
$d_1^{(r)},\ldots,d_{i-1}^{(r)}$, the mutable rows are those computed in
Appendix~\ref{app:triangular-mutation}, up to simultaneous reversal of
all arrows.  Proposition~\ref{prop:triangular-mutation-sequence} gives
\eqref{eq:positive-removal-row}.  By the block and boundary-identification
relations, every additional column incident with a mutated row lies on the
adjacent unmutated side; those
columns are retained in Appendix~\ref{app:triangular-mutation}.  Thus every
entry coming from an identified boundary is included in the calculation.
After the source or sink is deleted, principal restriction commutes with every
remaining mutation; the reverse sequence is therefore the inverse of the
forward sequence on the restricted matrix.

The odd half-diamond is obtained from the even one by
\[ (i,j)\longmapsto(\ell-i-j,j) \]
and global arrow reversal.  This sends the order $1,\ldots,N$ to
$N,\ldots,1$ and $\mathbf p_{N+1-i}^+(r)$ to $\mathbf p_i^-(r)$, proving
\eqref{eq:negative-removal-row}.
\end{proof}

\subsection{The diagonal reduction}

A flip of a disk diagonal in an $\ell$-triangulated quadrilateral is the
standard Fock--Goncharov mutation sequence
\cite[Section~10.3 and Proposition~10.1]{FockGoncharov}.  In
\cite[Step~2 in the proof of Theorem~6.22]{FeiKroneckerII}, this sequence is
applied at the transpose-paired diagonals $[k]$ and
$\overleftarrow{[k]}$ before the next diamond diagonal is removed.

\begin{lemma}[folding paired disk flips]
\label{lem:flip-support-separation}
After the exposed transpose-fixed diagonal has been deleted, the interiors
of the two quadrilaterals supporting the paired flip of \cite{FeiKroneckerII} have no arrows
between them.  Their standard flip sequences are transpose images and mutate no fixed
boundary vertex.  Hence the paired flip descends to the folded
matrix.
\end{lemma}

\begin{proof}
The Fock--Goncharov quiver is the sum of the oriented quivers of the small
triangles.  Once the common fixed diagonal is removed, no small triangle
contains an interior vertex from each quadrilateral.  Their interiors thus
form separated sets $L$ and $R=\tau L$; the shared boundary is fixed and
unmutated.  Apply Lemma~\ref{lem:separated-pair}.
\end{proof}

\begin{proposition}[folded diagonal reduction]
\label{prop:direct-folded-diagonal-reduction}
If $m\ge3$ is odd, then $B_{\ell,m}^{\mathrm{uf}}$ can be reduced by
mutations and source-sink deletions to the string-diagram exchange matrix
for the Cartan matrix $A_{\ell-1}$ associated with a triangulation of an
$(m+1)$-gon.
\end{proposition}

\begin{proof}
Remove the outer diagonal $m$ by
Lemma~\ref{lem:folded-half-diamond-removal}.  Inductively, suppose that the
matrix obtained at this point is the fold of the disk matrix in \cite{FeiKroneckerII}
after deleting diagonals
$m,m-1,\ldots,k+1$.  Apply the folded paired flip of
Lemma~\ref{lem:flip-support-separation}, remove the now-exposed diagonal $k$ by
Lemma~\ref{lem:folded-half-diamond-removal}, and undo the flip.  Reversing the preparatory sequences restores all entries between surviving
vertices, so the
induction hypothesis is preserved with diagonal $k$ deleted.

After all fixed diamond diagonals are removed, the disk decomposition in
\cite{FeiKroneckerII}
leaves two disconnected transpose-exchanged copies of the
$\ell$-triangulation of an $(m+1)$-gon
\cite[Step~3 in the proof of Theorem~6.22]{FeiKroneckerII}.  Folding leaves
one copy, whose mutable matrix is the stated string-diagram matrix.
\end{proof}

\begin{lemma}[compatibility as $m$ increases]
\label{lem:successive-principal-submatrix}
Let $\iota_m$ send every orbit label represented in a block $n\le m$ to the
same orbit label in $\mathscr V_{\ell,m+1}$.  Its restriction to
$\mathscr M_{\ell,m}$ is injective, and
\begin{equation}
 b^{(\ell,m+1)}_{\iota_m(u),\iota_m(v)}
 =b^{(\ell,m)}_{u,v}
 \qquad(u,v\in\mathscr M_{\ell,m}).
 \label{eq:successive-m-entry-equality}
\end{equation}
Consequently $B_{\ell,m}^{\mathrm{uf}}$ is the principal submatrix of
$B_{\ell,m+1}^{\mathrm{uf}}$ on $\iota_m(\mathscr M_{\ell,m})$.
\end{lemma}

\begin{proof}
The map is the identity on representatives from blocks $2,\ldots,m$.  The
only new equivalence relation is the boundary identification joining block $m$ to block
$m+1$, and it has one of the following two forms.

If $m$ is even, then $m+1$ is odd and
\begin{equation}
 (m+1;i,\ell-i)\sim(m;i,\ell-i)
 \qquad(1\le i\le\ell-1).
 \label{eq:successive-even-old-boundary}
\end{equation}
The labels on the right form the frozen outer boundary of
$\Sigma_{\ell,m}$; after the new block is attached they form an internal
identified boundary.  If $m$ is odd, then $m+1$ is even and
\begin{equation}
 (m+1;0,j)\sim(m;0,j)
 \qquad(1\le j\le\ell-1),
 \label{eq:successive-odd-old-boundary}
\end{equation}
again turning the old frozen outer boundary into an internal identified
boundary.  Thus no two labels that were mutable in $\Sigma_{\ell,m}$ are newly
identified, which proves injectivity on $\mathscr M_{\ell,m}$.

For the matrix entries, every relation supported wholly in blocks
$2,\ldots,m-1$ is literally unchanged.  In block $m$, the base-edge and
interior relations \eqref{eq:base-edge}--\eqref{eq:interior} are also
unchanged.  The only relations replaced on adjoining the new block are the
terminal-boundary relations: they become the even boundary-identification relations
\eqref{eq:even-gluing} in the first parity case and the odd boundary-identification relations
\eqref{eq:odd-gluing} in the second.  In either formula, every new factor is
either a variable from the new block or one of the old terminal labels in
\eqref{eq:successive-even-old-boundary} or
\eqref{eq:successive-odd-old-boundary}.  Hence no new exchange monomial
contains two old mutable labels, and no exponent relating two old mutable
labels changes.  This proves \eqref{eq:successive-m-entry-equality}.
Deleting the newly mutable former boundary labels and all labels belonging
only to block $m+1$ therefore leaves exactly $B_{\ell,m}^{\mathrm{uf}}$.
\end{proof}

\begin{theorem}\label{thm:reddening}
For every $\ell,m\ge2$, the seed $\Sigma_{\ell,m}$ admits a reddening
sequence.
\end{theorem}

\begin{proof}
Suppose first that $m$ is odd.  Proposition
\ref{prop:direct-folded-diagonal-reduction} reduces
$B_{\ell,m}^{\mathrm{uf}}$ to a string-diagram matrix.  The latter lies in
$\mathscr P'$; reading the reduction backwards shows that
$B_{\ell,m}^{\mathrm{uf}}\in\mathscr P'$.  It therefore has a reddening
sequence \cite[Theorem~4.8 and Corollary~4.9]{CaoString}.

If $m$ is even, then $m+1$ is odd and
$B_{\ell,m+1}^{\mathrm{uf}}$ has a reddening sequence.  By
Lemma~\ref{lem:successive-principal-submatrix},
$B_{\ell,m}^{\mathrm{uf}}$ is a principal submatrix of it.  Green-to-red
sequences pass to principal submatrices of skew-symmetrizable matrices
\cite[Lemma~4.9]{CaoLi}.
\end{proof}

\section{Theta bases and polyhedral formulas for symmetric-square plethysm}
\label{sec:polyhedral-plethysm}

We work over $\mathbb C$ in the theta-function arguments.  Lemma~\ref{lem:theta-characteristic-zero-descent} descends the resulting bases to the ground field fixed in
Section~\ref{sec:plethysm-highest-weight}.  Gross--Hacking--Keel--Kontsevich
construct theta functions from cluster scattering diagrams and prove a
corrected form of the Fock--Goncharov dual-basis conjecture
\cite[Introduction, Sections~7--8]{GHKK}.  Cheung--Magee--Mandel--Muller
prove valuative independence and theta reciprocity for seed data and obtain
theta bases for anticanonical partial compactifications
\cite{CMMM}.

\subsection{The folded seed datum and the open theta basis}

Invert the frozen variables and set
\[ \cU_{\ell,m}^{\circ} :=\cU_{\ell,m}[z_f^{-1}:f\in\mathscr F_{\ell,m}]. \]
Let $\mathcal A_{\ell,m}^{\circ}$ be the open cluster $A$-variety obtained
by inverting every frozen coordinate.  For each seed $t$ in the mutation
class, put
\[ \mathcal A_t :=\operatorname{Spec} \kk[z_{u;t}^{\pm1}:u\in\mathscr M_{\ell,m}] [z_f:f\in\mathscr F_{\ell,m}], \]
and glue these charts by the usual mutable cluster transformations.  Their
union inside the common cluster function field is denoted by
$\mathcal A_{\ell,m}$.

The initial chart identifies the theta lattice with
\[ M_{\ell,m}:=\mathbb Z^{\mathscr I_{\ell,m}}, \qquad M_{\ell,m,\mathbb R}:=M_{\ell,m}\otimes_{\mathbb Z}\mathbb R, \qquad N_{\ell,m}:=\operatorname{Hom}(M_{\ell,m},\mathbb Z). \]
Write $e_v$ for the standard basis of $M_{\ell,m}$ and $e_v^\vee$ for the
dual basis of $N_{\ell,m}$.  Let
\[ L_{\ell,m}:=\mathbb Z^{\mathscr M_{\ell,m}} \]
with basis $\epsilon_u$ indexed by mutable vertices.

\begin{proposition}[the folded seed datum]
\label{prop:theta-seed-datum}
Define homomorphisms
\begin{equation*}
 \begin{aligned}
 P:L_{\ell,m}&\longrightarrow M_{\ell,m},
 &P(\epsilon_u)&=\sum_{v\in\mathscr I_{\ell,m}}b_{uv}e_v,\\
 Q:L_{\ell,m}&\longrightarrow N_{\ell,m},
 &Q(\epsilon_u)&=e_u^\vee.
 \end{aligned}
\end{equation*}
Let $D=\operatorname{diag}(d_u)$, where $d_u$ is the transpose-orbit size
from Proposition~\ref{prop:folded-B}, and put
\begin{equation}
 \begin{aligned}
 N_{\ell,m}^{\bullet}
 &:=N_{\ell,m}
   +\sum_{u\in\mathscr M_{\ell,m}}
     \mathbb Z\frac{e_u^\vee}{d_u},\\
 M_{\ell,m}^{\bullet}
 &:=\bigoplus_{u\in\mathscr M_{\ell,m}}\mathbb Z\,d_ue_u
   \ \oplus\!
   \bigoplus_{f\in\mathscr F_{\ell,m}}\mathbb Z\,e_f,\\
 Q^{\bullet}&:=QD^{-1},
 &P^{\bullet}&:=PD.
 \end{aligned}
 \label{eq:bullet-lattices-maps}
\end{equation}
Then $(P,Q)$ is a skew-symmetrizable seed datum in the sense of
\cite[Section~4.1]{CMMM}.  Its $A$-type scattering diagram is the cluster
scattering diagram of $\Sigma_{\ell,m}$, and its chiral dual, in the terminology of \cite[Section~4.5]{CMMM}, is the seed
datum
\[ (Q^{\bullet},P^{\bullet}) \quad\text{on the dual lattices}\quad (N_{\ell,m}^{\bullet},M_{\ell,m}^{\bullet}). \]
If
\[ \mathbf B:=Q^{\mathsf T}P=(B_{\ell,m}^{\mathrm{uf}})^{\mathsf T}, \qquad \mathbf B^{\bullet}:=(Q^{\bullet})^{\mathsf T}P, \]
then
\begin{equation}
 \mathbf B=D\mathbf B^{\bullet},
 \qquad
 \mathbf B^{\bullet}_{uv}=\frac{b_{vu}}{d_u}
 =-\frac{b_{uv}}{d_v}.
 \label{eq:CMMM-skew-matrix}
\end{equation}
In particular, $\mathbf B^{\bullet}$ is skew-symmetric over $\mathbb Q$.
\end{proposition}

\begin{proof}
The lattices in \eqref{eq:bullet-lattices-maps} are mutually dual.  The map
$Q^{\bullet}$ is integral by construction.  For a mutable coordinate $v$,
the coefficient of $e_v$ in $P^{\bullet}(\epsilon_u)$ is
\[ d_ub_{uv}=-d_vb_{vu}, \]
which is divisible by $d_v$; frozen coordinates impose no further
divisibility condition.  Hence $P^{\bullet}$ takes values in
$M_{\ell,m}^{\bullet}$.  Equation \eqref{eq:CMMM-skew-matrix} follows from
\eqref{eq:symmetrizer}, and proves the seed-datum axioms.  With the mutable
indices ordered first, $P$ is the transpose of the extended exchange matrix
and $Q$ is the block inclusion $\binom{I}{0}$.  Thus
\cite[Example~4.2]{CMMM} identifies $\mathcal D(P,Q)$ with the $A$-type
cluster scattering diagram for $\Sigma_{\ell,m}$.  The description of the
chiral dual is \cite[Section~4.5]{CMMM}.
\end{proof}


\begin{proposition}\label{prop:open-theta-basis}
For every $\ell,m\ge2$,
\[ \{\vartheta_g:g\in M_{\ell,m}\} \]
is a basis of $\cU_{\ell,m}^{\circ}$.
\end{proposition}

\begin{proof}
The terminal $c$-vectors of the reddening sequence in
Theorem~\ref{thm:reddening} are nonpositive.  By
\cite[Lemma~5.12]{GHKK}, they are the inward normals to the terminal cluster
chamber.  That chamber therefore meets the interior of the initial negative
chamber, so \cite[Lemma~8.29]{GHKK} gives a large cluster complex.  This is
the argument for $(2)\Rightarrow(3)$ in \cite[Corollary~8.30]{GHKK}; it uses
only the terminal sign condition.  Proposition~\ref{prop:folded-B} verifies
the convexity condition in \cite[Theorem~0.3(7)]{GHKK}.  Hence
\cite[Proposition~8.28 and Definition~0.6]{GHKK} identifies the middle,
upper, and canonical algebras.  The canonical basis is indexed by all
integral tropical points of the dual, which the initial seed identifies
with $M_{\ell,m}$.
\end{proof}

\begin{remark}[relation with Fock--Goncharov]
\label{rem:Fock-Goncharov-relation}
Two parts of the construction come from the work of Fock and Goncharov.  Their
local-system construction supplies the $\ell$-triangulated disk quivers and
the flip sequences used in Section~\ref{sec:reddening}
\cite[Section~10]{FockGoncharov}.  Their cluster-ensemble formalism
abstracts these examples into a pair of positive spaces of $A$- and
$X$-type, related by a map $p$, and formulates the duality conjectures in
terms of integral tropical points
\cite[Sections~1--4]{FockGoncharovEnsembles}.  The present folded invariant
variety is not asserted to be a moduli space of local systems.  What is used
is the disk mutation geometry before folding and the $A/X$ duality
formalism after folding.  The orbit-size symmetrizer makes the chiral dual
of Proposition~\ref{prop:theta-seed-datum}, rather than a naive transpose,
the appropriate mirror seed datum.  The cone $\mathscr C_{\ell,m}$ below is
the set of tropical parameters whose theta functions extend across the
frozen boundary.
\end{remark}

\subsection{The frozen boundary and its theta basis}

\begin{proposition}\label{prop:frozen-partial-compactification}
The scheme $\mathcal A_{\ell,m}$ is an integral partial compactification of
$\mathcal A_{\ell,m}^{\circ}$ with the following properties.
\begin{enumerate}
 \item Its global functions are the upper cluster algebra with polynomial
       frozen coefficients, while localization at all frozen variables
       gives the open upper cluster algebra:
       \begin{equation}
       \Gamma(\mathcal A_{\ell,m},\mathcal O)=\cU_{\ell,m},
       \qquad
       \Gamma(\mathcal A_{\ell,m}^{\circ},\mathcal O)
       =\cU_{\ell,m}^{\circ}.
       \label{eq:cluster-global-sections}
       \end{equation}
 \item For each frozen index $f$, the equations $z_f=0$ in the seed charts
       glue to an irreducible effective Cartier divisor $D_f$, with
       $\operatorname{ord}_{D_f}(z_f)=1$.  The complement of the open
       cluster variety is
       \begin{equation}
       \mathcal A_{\ell,m}\setminus\mathcal A_{\ell,m}^{\circ}
       =\bigcup_{f\in\mathscr F_{\ell,m}}D_f,
       \label{eq:frozen-boundary-union}
       \end{equation}
       and it has no other codimension-one components.
 \item The divisor $D=\sum_fD_f$ is anticanonical.  More precisely, the
       logarithmic volume forms
       \begin{equation}
       \Omega_t
       =\bigwedge_{u\in\mathscr M_{\ell,m}}d\log z_{u;t}
        \wedge
        \bigwedge_{f\in\mathscr F_{\ell,m}}d\log z_f
       \label{eq:cluster-volume-form}
       \end{equation}
       glue up to sign, and
       $\operatorname{div}(\Omega)=-\sum_fD_f$.
\end{enumerate}
Hence $(\mathcal A_{\ell,m},D)$ satisfies the frozen anticanonical
partial-compactification hypotheses used in
\cite[Corollary~1.2]{CMMM}.
\end{proposition}

\begin{proof}
Every chart ring is a Laurent polynomial ring in the mutable coordinates
and an ordinary polynomial ring in the common frozen coordinates.  The
charts are dense affine opens of one function field and the gluing is
separated, so a global regular function is exactly an element regular in
every chart.  Intersecting their coordinate rings gives the defining
upper-cluster intersection with polynomial frozen coefficients.  Inverting
all $z_f$ gives the corresponding open intersection, proving (1).

In any chart, $(z_f)$ is a height-one prime ideal and $z_f$ is a local
parameter.  Frozen coordinates are unchanged by mutable mutation.  On the
overlap for a mutation at $u$, exactly one exchange monomial can contain
$z_f$, because the single matrix entry $b_{uf}$ has one sign; after setting
$z_f=0$, the other monomial is generically nonzero.  Hence the two chart
pieces of $(z_f=0)$ have a common dense open subset and glue to the closure
of the irreducible initial-chart divisor $(z_f=0)$.  This proves that $D_f$
is an irreducible prime Cartier divisor, while distinct frozen indices give
distinct generic points in the initial chart.  Localizing by $\prod_fz_f$
removes exactly their union.  Thus every codimension-one component of the
complement is one of the $D_f$, proving (2).

A direct Jacobian calculation for one cluster mutation gives
$\Omega_{\mu_u(t)}=-\Omega_t$; hence the forms in
\eqref{eq:cluster-volume-form} glue up to sign.  At the generic point of
$D_f$, the factor $d\log z_f=dz_f/z_f$ has a simple pole and every other
coordinate in this chart is a unit.  Therefore
$\operatorname{ord}_{D_f}(\Omega)=-1$.  By
\eqref{eq:frozen-boundary-union} there is no other boundary divisor on
which a pole can occur, which proves (3).
\end{proof}

\begin{proposition}\label{prop:primitive-frozen-valuations}
For $f\in\mathscr F_{\ell,m}$, the divisorial valuation of $D_f$ is the
primitive vector
\[ n_f=e_f^\vee\in N_{\ell,m}\subset N_{\ell,m}^{\bullet}. \]
In the coordinate lattice of every seed obtained by mutable mutations,
$n_f$ still has value $1$ on the frozen coordinate $z_f$, value $0$ on all
other frozen coordinates, and value $0$ on every mutable cluster
coordinate.
\end{proposition}

\begin{proof}
In the initial torus chart,
$\operatorname{val}_{D_f}(z^g)=g_f=\langle g,e_f^\vee\rangle$, 
so the valuation vector is exactly $e_f^\vee$, not a nontrivial multiple.
It is primitive because $\operatorname{val}_{D_f}(z_f)=1$.

Assume inductively that all mutable cluster coordinates obtained so far
have $D_f$-valuation zero.  Mutation at $u$
has exchange relation
\[ z_uz_u' =\prod_vz_v^{[b_{uv}]_+} +\prod_vz_v^{[-b_{uv}]_+}. \]
The two monomials have $D_f$-valuations $[b_{uf}]_+$ and
$[-b_{uf}]_+$, so at least one has valuation zero.  Their sum is nonzero at
the generic point of $D_f$, and hence has valuation zero.  Since
$\operatorname{val}_{D_f}(z_u)=0$, also
$\operatorname{val}_{D_f}(z_u')=0$.  Induction proves the claim.
\end{proof}

For $f\in\mathscr F_{\ell,m}$ define the boundary valuation function
\[ \nu_f:M_{\ell,m,\mathbb R}\longrightarrow\mathbb R, \qquad \nu_f(g):=\operatorname{val}_{D_f}(\vartheta_g), \]
and set
\begin{equation}
 \mathscr C_{\ell,m}
 :=\left\{g\in M_{\ell,m,\mathbb R}:
   \nu_f(g)\ge0\text{ for every }f\in\mathscr F_{\ell,m}\right\}.
 \label{eq:theta-cone}
\end{equation}
By valuative independence and
\cite[Theorem~1.1 and Corollary~1.2]{CMMM}, the functions $\nu_f$ extend to
convex integral piecewise-linear functions and \eqref{eq:theta-cone} is a
polyhedron.  Each $\nu_f$ is an infimum of integral linear forms, hence is
positively homogeneous.  Because the right-hand sides in
\eqref{eq:theta-cone} are zero, $\mathscr C_{\ell,m}$ is therefore a rational
polyhedral cone.  After clearing denominators, choose an integral matrix
$H_{\ell,m}$ such that
\[ \mathscr C_{\ell,m} =\{g\in M_{\ell,m,\mathbb R}:H_{\ell,m}g\ge0\}. \]

For $n\in N_{\ell,m}^{\bullet}$ write
\[ \vartheta_n^{\chi}:=\Theta_{n,+}^{(Q^{\bullet},P^{\bullet})}. \]
Theta reciprocity \cite[Claim~4.16 and Theorem~5.19]{CMMM}, together with
Proposition~\ref{prop:primitive-frozen-valuations}, gives
\[ \nu_f(g)=\operatorname{val}_{g}(\vartheta_{n_f}^{\chi}) \]
whenever both sides are finite.  Thus, if $\vartheta_{n_f}^{\chi}=\sum_{a\in A_f}c_{f,a}X^a \ (c_{f,a}>0)$,
is a finite Laurent expansion in the initial chiral-dual chart, then
\[ \nu_f(g)=\min_{a\in A_f}\langle a,g\rangle. \]

\begin{remark}[choice of the chiral dual]
\label{rem:chiral-versus-Langlands}
The standard Langlands-dual pair is $(Q,-P)$.  The corresponding integral
seed-datum version of theta reciprocity requires both $D$ and
$\mathbf B^{\bullet}$ to be integral, and it uses the theta function
expanded from the \emph{negative} chamber:
\begin{equation}
 \operatorname{val}_{n}(\Theta_{g,+}^{(P,Q)})
 =\operatorname{val}_{g}(\Theta_{n,-}^{(Q,-P)}).
 \label{eq:Langlands-negative-chamber}
\end{equation}
By \cite[Corollary~4.20]{CMMM}, in the folded seed $D$ is integral,
but integrality of $\mathbf B^{\bullet}$ does not follow from
skew-symmetrizability and is not assumed.  Consequently all reciprocity and
boundary-potential formulas in this paper use the chiral dual
$(Q^{\bullet},P^{\bullet})$.  Formula
\eqref{eq:Langlands-negative-chamber} is available in any individual case
where $\mathbf B^{\bullet}$ is integral, but the negative-chamber
convention must then be retained.
\end{remark}

\begin{theorem}[theta basis of the partial compactification]
\label{thm:theta-basis}
For every $\ell,m\ge2$,
\[ \left\{\vartheta_g: g\in\mathscr C_{\ell,m}\cap M_{\ell,m}\right\} \]
is a basis of $\cR_{\ell,m}$.
\end{theorem}

\begin{proof}
Proposition~\ref{prop:open-theta-basis} gives a theta basis on the open
cluster variety.  Corollary~1.2 of \cite{CMMM} and valuative independence
say that the members regular along every $D_f$ form a basis of the partial
compactification.  By definition, regularity is exactly the condition
$g\in\mathscr C_{\ell,m}$.  Finally,
\eqref{eq:cluster-global-sections} and Theorem~\ref{thm:main-equality}
identify the resulting algebra with $\cR_{\ell,m}$.
\end{proof}

\subsection{Weight fibers and plethysm}

Define the initial weight map, using the weights of
Proposition~\ref{prop:weights-count}, by
\[ \mathsf W_{\ell,m}:M_{\ell,m} \longrightarrow X^*(T_V)\oplus\mathbb Z^m, \qquad e_v\longmapsto\bigl(\wt_V(z_v),\deg_W(z_v)\bigr). \]
Exchange homogeneity implies
\[ \mathsf W_{\ell,m}(b_{u\bullet})=0 \qquad(u\in\mathscr M_{\ell,m}). \]
Thus $\widehat y_u=\prod_vz_v^{b_{uv}}$ has weight zero, and a pointed
theta expansion yields
\begin{equation}
 \wt(\vartheta_g)=\mathsf W_{\ell,m}g.
 \label{eq:theta-weight}
\end{equation}

For a dominant polynomial $\GL_\ell$-weight $\lambda$ and a weak composition
$\alpha\in\mathbb Z_{\ge0}^m$, set
\[ \mathscr P_{\ell,m}(\lambda;\alpha) :=\left\{g\in M_{\ell,m,\mathbb R}: H_{\ell,m}g\ge0, \quad \mathsf W_{\ell,m}g=(\lambda;\alpha) \right\}. \]

\begin{corollary}\label{cor:polyhedral-b}
For every $\ell,m\ge2$, we have
\begin{equation}
 b_{\alpha,(2)}^\lambda
 =\left|\mathscr P_{\ell,m}(\lambda;\alpha)\cap M_{\ell,m}\right|.
 \label{eq:polyhedral-b}
\end{equation}
\end{corollary}

\begin{proof}
By Theorem~\ref{thm:theta-basis} and \eqref{eq:theta-weight}, these
lattice points index a basis of the $(\lambda;\alpha)$-weight space of
$\cR_{\ell,m}$.  Its dimension is $b_{\alpha,(2)}^\lambda$ by
\eqref{eq:b-as-multiplicity}.  The set is finite because this
multihomogeneous component of the polynomial invariant ring is
finite-dimensional.
\end{proof}

\begin{theorem}[polyhedral formula for symmetric-square plethysm]
\label{thm:polyhedral-plethysm}
Let $\ell,m\ge2$, and let $\lambda,\mu$ be partitions with
$\htp(\lambda)\le\ell$ and $\htp(\mu)\le m$.  Then
\[  a_{\mu,(2)}^\lambda =\sum_{\sigma\in\frakS_m}\operatorname{sgn}(\sigma) \left| \mathscr P_{\ell,m} \bigl(\lambda;\alpha^\sigma(\mu)\bigr) \cap M_{\ell,m} \right|. \]
Here $\alpha^\sigma(\mu)_i=\mu_i-i+\sigma(i) \ (1\le i\le m)$,
and a summand is zero if $\alpha^\sigma(\mu)$ has a negative component.
\end{theorem}

\begin{proof}
Substitute \eqref{eq:polyhedral-b} into the alternating formula of
Theorem~\ref{thm:alternating-plethysm}; Theorem~\ref{thm:main-equality}
identifies the cluster algebra used in the count with the invariant ring.
\end{proof}

\subsection{Optimized boundary seeds for odd rank and an even number of matrices}

The valuation-theoretic statements above hold for all $(\ell,m)$.  When
$\ell$ is odd and $m$ is even, each boundary valuation is optimized in a
seed obtained by mutation.

Let $\Sigma'$ be a seed in the mutation class, with seed map $P'$.  For a
frozen index $f$, Proposition~\ref{prop:primitive-frozen-valuations} and
$P'^{\bullet}=P'D$ give
\[ \langle n_f,P'(\epsilon_u)\rangle=b'_{uf}. \]
Hence $\Sigma'$ is optimized for $n_f$ precisely when $b'_{uf}\ge0$ for
every mutable $u$.  In this case $\vartheta_{n_f}^{\chi}$ restricts to
$X^{n_f}$ on the corresponding chiral-dual torus by
\cite[Definition~9.1 and Lemma~9.3]{GHKK}.

\begin{proposition}\label{prop:optimized-boundary-seeds}
Assume that $\ell$ is odd and $m$ is even.  Every frozen boundary divisor
admits an optimized seed.  The seed may be chosen so that the corresponding
frozen column has a single nonzero mutable entry.
\end{proposition}

\begin{proof}
For a determinant divisor $f=\partial_r$, apply the preparatory mutations of
\cite{FeiKroneckerII} in the odd diamonds and mutate along the $r$-th
diamond diagonal toward $\partial_r$.  Lemma~\ref{lem:folded-preparatory-twists} folds the twists on
the principal matrix and the chosen frozen column; the remaining diagonal
vertices are transpose-fixed.  The hive-row calculation in \cite{FeiKroneckerII} leaves
$\partial_r$ incident with one mutable vertex, with the arrow directed into
$\partial_r$ \cite[Lemmas~6.12--6.13]{FeiKroneckerII}.

Let $f=[m;i,\ell-i]$ be an outer frozen divisor and choose the terminal
vertex $v=(_{i,\ell-i}^{m})$ in the unfolded seed.  For even $m$, the other
terminal vertex of the full boundary path in \cite{FeiKroneckerII} is
$v^{\star}=(_{\ell-i,i}^{m})$.
These terminal vertices lie in the same half of the terminal disk;
transposition sends them to the reflected half.  Their folded orbits coincide only when
$i=\ell-i$, which is impossible for odd $\ell$.  Pair the path mutation
sequence of \cite{FeiKroneckerII}, oriented from $v^{\star}$ to $v$, with its transpose.
Proposition~\ref{prop:folded-even-boundary-path} shows that the sequence folds
on the principal matrix and the chosen frozen column.  The terminal column
is
$B^{\mathbf q_f}_{\bullet f}=c e_u$ ($c>0$),
for one mutable index $u$, so the resulting seed is optimized.
\end{proof}

For odd $\ell$ and even $m$, put
\[ W_{\partial}^{\chi} :=\sum_{f\in\mathscr F_{\ell,m}}\vartheta_{n_f}^{\chi}. \]

\begin{corollary}\label{cor:boundary-cone-optimized-seeds}
Assume that $\ell$ is odd and $m$ is even.  For each frozen index $f$,
choose an optimizing mutation sequence and write the pullback of $X^{n_f}$ to
the initial chiral-dual chart as
\begin{equation}
 \vartheta_{n_f}^{\chi}=\sum_{a\in A_f}c_{f,a}X^a,
 \qquad c_{f,a}>0.
 \label{eq:boundary-theta-Laurent}
\end{equation}
Then
\begin{equation}
 \mathscr C_{\ell,m}
 =\left\{g:\langle a,g\rangle\ge0
   \text{ for every }f\in\mathscr F_{\ell,m},\ a\in A_f\right\}.
 \label{eq:optimized-boundary-cone}
\end{equation}
Clearing denominators and removing redundant forms gives an integral
inequality matrix $H_{\ell,m}$.
\end{corollary}

\begin{proof}
By the optimized-seed criterion above, the theta function is a monomial in
the optimized seed.  Pulling it back through the inverse chiral-dual
mutation sequence gives \eqref{eq:boundary-theta-Laurent}.  Theta reciprocity yields
\[ \nu_f(g)=\min_{a\in A_f}\langle a,g\rangle, \]
so regularity along every frozen divisor is equivalent to
\eqref{eq:optimized-boundary-cone}.  Equivalently,
\[ \mathscr C_{\ell,m} =\{g:(W_{\partial}^{\chi})^{\operatorname{trop}}(g)\ge0\}. \]
\end{proof}

\subsection{Successive matrix models as degree-zero faces}
\label{subsec:successive-matrix-faces}

Let $m\ge3$.  For $1\le s\le m$, let
\[ d_s:=\operatorname{pr}_s\circ\mathsf W_{\ell,m}:M_{\ell,m,\mathbb R}\longrightarrow\mathbb R, \]
where $\operatorname{pr}_s$ is projection to the $s$th $T_W$-degree.  For
an integral theta parameter $g$, equation~\eqref{eq:theta-weight}
identifies $d_s(g)$ with the degree of $\vartheta_g$ in the entries of
$C_s$.  If $2\le r<m$, put
\[ \mathscr C_{\ell,r}^{[m]}:=\mathscr C_{\ell,m}\cap\bigcap_{s=r+1}^m\ker(d_s), \qquad \delta_{r,m}:=\sum_{s=r+1}^m d_s. \]
The superscript $[m]$ records the ambient cone from which the model is
inherited; set $\mathscr C_{\ell,m}^{[m]}:=\mathscr C_{\ell,m}$.

\begin{proposition}\label{prop:successive-matrix-face}
For every $\ell\ge2$, $m\ge3$, and $2\le r<m$, each $d_s$ is nonnegative
on $\mathscr C_{\ell,m}$, and
\[ \mathscr C_{\ell,r}^{[m]}=\mathscr C_{\ell,m}\cap\ker(\delta_{r,m}) \]
is a proper exposed rational polyhedral face of $\mathscr C_{\ell,m}$.  Under the
pullback $\cR_{\ell,r}\hookrightarrow\cR_{\ell,m}$ induced by the
projection $(C_1,\ldots,C_m)\mapsto(C_1,\ldots,C_r)$, the family
\[ \{\vartheta_g:g\in\mathscr C_{\ell,r}^{[m]}\cap M_{\ell,m}\} \]
is a basis of $\cR_{\ell,r}$.  Thus, for
$\alpha\in\mathbb Z_{\ge0}^r$ and a dominant polynomial
$\GL_\ell$-weight $\lambda$,
\[
 b_{\alpha,(2)}^\lambda
 =\left|\left\{g\in\mathscr C_{\ell,r}^{[m]}\cap M_{\ell,m}:
 \mathsf W_{\ell,m}g=
 (\lambda;\alpha,\underbrace{0,\ldots,0}_{m-r})\right\}\right|.
\]
If $H_{\ell,m}$ is an inequality matrix for $\mathscr C_{\ell,m}$, this
model is obtained by imposing
\[ H_{\ell,m}g\ge0,\qquad d_{r+1}(g)=\cdots=d_m(g)=0, \]
and restricting the weight map to the first $r$ matrix degrees.
\end{proposition}

\begin{proof}
Let $g\in\mathscr C_{\ell,m}\cap M_{\ell,m}$.  By
Theorem~\ref{thm:theta-basis}, $\vartheta_g$ is a nonzero polynomial
invariant in $\cR_{\ell,m}$, and
\eqref{eq:theta-weight} makes it homogeneous of multidegree
$(d_1(g),\ldots,d_m(g))$.  Hence $d_s(g)\ge0$ for every $s$.  A rational
polyhedral cone is generated by lattice vectors, so each $d_s$ is
nonnegative on all of $\mathscr C_{\ell,m}$.  It follows that
$\delta_{r,m}\ge0$ on the cone and that its zero locus there is precisely
$\mathscr C_{\ell,r}^{[m]}$.  This proves the face assertion.

The simultaneous degree-zero part of the polynomial ring on $m$ symmetric
matrices, with respect to $C_{r+1},\ldots,C_m$, is the polynomial ring on
the first $r$ matrices.  The congruence action preserves the full
multigrading, and therefore
\[ (\cR_{\ell,m})_{d_{r+1}=\cdots=d_m=0}=\cR_{\ell,r}. \]
The theta basis is homogeneous, so the members with these degrees zero form
a basis of this subalgebra.  Since $\Delta_m$ has positive
$\delta_{r,m}$-degree, some homogeneous theta basis element lies outside the
face; hence the face is proper.  The counting formula and the description by
restricted inequalities and weights follow.
\end{proof}

For fixed $m$, these faces form the chain
\[ \mathscr C_{\ell,2}^{[m]}\subset\mathscr C_{\ell,3}^{[m]}\subset\cdots\subset\mathscr C_{\ell,m-1}^{[m]}\subset\mathscr C_{\ell,m}. \]
For $2\le r<m$,
\[ \mathscr C_{\ell,r}^{[m]}=\mathscr C_{\ell,r+1}^{[m]}\cap\ker(d_{r+1}), \]
so each inclusion is an exposed face.  Taking $m=2n$ and $r=2n-1$ gives a
polyhedral theta model for $\cR_{\ell,2n-1}$ as a face of
$\mathscr C_{\ell,2n}$.  If $\ell$ is odd,
Corollary~\ref{cor:boundary-cone-optimized-seeds} supplies the explicit
matrix $H_{\ell,2n}$, which may then be restricted by $d_{2n}=0$.

\subsection{Descent to the ground field}

\begin{lemma}\label{lem:theta-characteristic-zero-descent}
Let $A_{\mathbb Q}$ be the polynomial ring in the matrix entries defining
$X_{\ell,m}$, and put
\[ R_{\mathbb Q}:=A_{\mathbb Q}^{U_{\ell,\mathbb Q}}. \]
For every algebraically closed field $K$ of characteristic zero, scalar
extension induces
\begin{equation}
 R_{\mathbb Q}\otimes_{\mathbb Q}K
 \simeq K[X_{\ell,m}]^{U_{\ell,K}}.
 \label{eq:invariant-base-change}
\end{equation}
The regular finite theta functions used in this section belong to the common
$\mathbb Q$-form of the cluster fraction field.  If they form a theta basis
over $\mathbb C$, then the same indexed family is a basis over every such
$K$.
\end{lemma}

\begin{proof}
The group $U_{\ell,\mathbb Q}$ is generated by its simple-root
one-parameter subgroups.  On every fixed polynomial multidegree, invariants
are the common kernel of finitely many $\mathbb Q$-linear locally nilpotent
derivations between finite-dimensional $\mathbb Q$-spaces.  Flat scalar
extension preserves this kernel, proving
\eqref{eq:invariant-base-change} degree by degree.

The chamber determinants and exchange relations are defined over
$\mathbb Z$.  The scattering functions and broken-line multiplication rules
are integral; in particular, every regular theta function that is finite in
an initial cluster chart has a Laurent expansion with integral
coefficients.  Such an element lies in the rational cluster fraction field.
If it is regular after extension to $\mathbb C$, its Laurent denominators
cancel already over $\mathbb Q$; equivalently,
\[ A_{\mathbb Q}=A_{\mathbb C}\cap\mathbb Q(X_{\ell,m}) \quad\text{inside }\mathbb C(X_{\ell,m}). \]
The same common-kernel description then puts it in $R_{\mathbb Q}$.

Fix a multidegree $\gamma$.  A theta family that is a basis of
$(R_{\mathbb Q}\otimes\mathbb C)_\gamma$ is linearly independent over
$\mathbb Q$.  By \eqref{eq:invariant-base-change}, its cardinality equals
$\dim_{\mathbb Q}(R_{\mathbb Q})_\gamma$, so it is already a
$\mathbb Q$-basis.  Tensoring with $K$ proves the claim.  The same argument
applies to the boundary-regular subfamily, because divisorial regularity is
the field-independent system of inequalities
$\operatorname{val}_{D_f}(\vartheta_g)\ge0$.
\end{proof}


\medskip
\noindent\emph{Software.}
The Python package distributed with the paper provides the routine
\texttt{bosonic\_model}.  Its main interface is
\[ \texttt{B, W, DT, H = bosonic\_model(l,m)}. \]
On input $(\ell,m)$, the first routine returns the extended exchange matrix $B_{\ell,m}$, the
variable-by-weight matrix $\mathsf W_{\ell,m}$, a certified reddening
sequence, and an integral inequality matrix $H_{\ell,m}$, with conventions
\[ B_{\ell,m}\mathsf W_{\ell,m}=0,\qquad \mathscr C_{\ell,m}=\{g:H_{\ell,m}g\geq0\}. \]
Mutable variables precede frozen variables, and mutation indices are
zero-based.  The routine \texttt{bosonic\_model\_data} also records the
variable labels, symmetrizer, weight coordinates, source model, and any
remaining face equations.  For odd $m$, the program obtains
$H_{\ell,m}$ from the even source model $(\ell,m+1)$ by Proposition~\ref{prop:successive-matrix-face}.  

\appendix

\section{Transfer from the flagged Kronecker model}
\label{app:matrix-transfer}

We prove Proposition~\ref{prop:matrix-transfer} and record the dictionary
between the lifted semi-invariants of \cite{FeiKroneckerII} and the matrix
chamber functions used
in Sections~\ref{sec:matrix-invariants} and \ref{sec:folding}.  The proof
proceeds through the flag-arm quotient, the Schofield presentations, and the
minimal lift.

\subsection{The flagged representation space}

Let $\mathsf K_{\ell,m}$ be the quiver with vertex set
\[ \{-\ell,\ldots,-1,1,\ldots,\ell\}. \]
It consists of two oppositely oriented type-$A$ arms joined by $m$ arrows
from $-\ell$ to $\ell$, as shown in Figure~\ref{fig:flagged-Klm}.

\begin{figure}[ht]
\centering
\begin{tikzpicture}[
  >=Latex,
  qvertex/.style={circle,draw,fill=white,minimum size=2.7mm,inner sep=0pt},
  qlabel/.style={font=\small},
  qdim/.style={font=\scriptsize},
  every path/.style={line width=.45pt}
]
\node[qvertex] (n1) at (0,0) {};
\node[qvertex] (n2) at (1.15,0) {};
\node (ndots) at (2.25,0) {$\cdots$};
\node[qvertex] (nlm) at (3.55,0) {};
\node[qvertex] (nl) at (5.05,0) {};
\node[qvertex] (pl) at (7.55,0) {};
\node[qvertex] (plm) at (9.05,0) {};
\node (pdots) at (10.35,0) {$\cdots$};
\node[qvertex] (p2) at (11.45,0) {};
\node[qvertex] (p1) at (12.60,0) {};

\node[qlabel] at (0,.38) {$-1$};
\node[qlabel] at (1.15,.38) {$-2$};
\node[qlabel] at (3.55,.38) {$-(\ell-1)$};
\node[qlabel] at (5.05,.38) {$-\ell$};
\node[qlabel] at (7.55,.38) {$\ell$};
\node[qlabel] at (9.05,.38) {$\ell-1$};
\node[qlabel] at (11.45,.38) {$2$};
\node[qlabel] at (12.60,.38) {$1$};

\draw[->] (n1) -- (n2);
\draw[->] (n2) -- (ndots);
\draw[->] (ndots) -- (nlm);
\draw[->] (nlm) -- (nl);
\draw[->] (pl) -- (plm);
\draw[->] (plm) -- (pdots);
\draw[->] (pdots) -- (p2);
\draw[->] (p2) -- (p1);

\draw[->] (nl) to[bend left=27] node[above] {$a_1$} (pl);
\draw[->] (nl) -- node[above] {$\cdots$} (pl);
\draw[->] (nl) to[bend right=27] node[below] {$a_m$} (pl);

\node[qdim,anchor=east] at (-.35,-.52) {$\boldsymbol\beta_\ell:$};
\node[qdim] at (0,-.52) {$1$};
\node[qdim] at (1.15,-.52) {$2$};
\node[qdim] at (2.25,-.52) {$\cdots$};
\node[qdim] at (3.55,-.52) {$\ell-1$};
\node[qdim] at (5.05,-.52) {$\ell$};
\node[qdim] at (7.55,-.52) {$\ell$};
\node[qdim] at (9.05,-.52) {$\ell-1$};
\node[qdim] at (10.35,-.52) {$\cdots$};
\node[qdim] at (11.45,-.52) {$2$};
\node[qdim] at (12.60,-.52) {$1$};
\end{tikzpicture}
\caption{The flagged $m$-arrow Kronecker quiver $\mathsf K_{\ell,m}$ and
its standard dimension vector.  The central bundle consists of the arrows
$a_r:-\ell\to\ell$ for $1\le r\le m$.}
\label{fig:flagged-Klm}
\end{figure}

Thus $\boldsymbol\beta_\ell(\pm i)=i$.  Write $V_{\pm i}=\kk^i$.  The
representation space decomposes as
\[ \Rep_{\boldsymbol\beta_\ell}(\mathsf K_{\ell,m}) =\Rep_{\boldsymbol\beta_\ell}(A_\ell) \times \Mat_\ell^m \times \Rep_{\boldsymbol\beta_\ell}(A_\ell^\vee), \]
where the middle factor records the central maps $A_r=M(a_r)$.  The
semi-invariant ring is
\[ \SI_{\boldsymbol\beta_\ell}(\mathsf K_{\ell,m}) :=\kk[\Rep_{\boldsymbol\beta_\ell}(\mathsf K_{\ell,m})] ^{\prod_v\SL(V_v)}. \]

Let $G=\SL_\ell$.  Quotienting the internal vertices of the negative and
positive arms gives the affine closures
\[ \mathcal A^-:=\operatorname{Spec}\kk[G]^{U^-}, \qquad \mathcal A^+:=\operatorname{Spec}\kk[G]^U \]
of $U^-\backslash G$ and $G/U$, respectively.  The first space carries a
right $G$-action and the second a left $G$-action.

The first fundamental theorem for the two type-$A$ flag arms identifies
the quotient by the internal base-change groups with
$\mathcal A^-\times\Mat_\ell^m\times\mathcal A^+$; see
\cite{Grosshans,DerksenWeyman}.  The two terminal copies of $G$ act by
\[ (F^-,A_1,\ldots,A_m,F^+) \longmapsto (F^-g^{-1},\,gA_1h^{-1},\ldots,gA_mh^{-1},\,hF^+). \]

Taking invariants under the two terminal groups gives
\[ \SI_{\boldsymbol\beta_\ell}(\mathsf K_{\ell,m}) \cong \kk[\mathcal A^-\times\Mat_\ell^m\times\mathcal A^+]^{G\times G}. \]

\subsection{The transfer isomorphism}

Choose the standard points $F^-\in\mathcal A^-$ and
$F^+\in\mathcal A^+$ represented by the standard complete flags.  Their
stabilizers are $U^-$ and $U$.  Evaluation induces
\begin{equation}
 \operatorname{ev}_{F^-,F^+}:
 \kk[\mathcal A^-\times\Mat_\ell^m\times\mathcal A^+]^{G\times G}
 \longrightarrow
 \kk[\Mat_\ell^m]^{U^-\times U}.
 \label{eq:app-evaluation}
\end{equation}

\begin{proposition}\label{prop:app-transfer}
The map \eqref{eq:app-evaluation} is a multigraded algebra isomorphism.
Consequently,
\begin{equation}
 \SI_{\boldsymbol\beta_\ell}(\mathsf K_{\ell,m})
 \xrightarrow{\sim} \kk[\Mat_\ell^m]^{U^-\times U}
 \xrightarrow{\sim} \cT_{\ell,m}.
 \label{eq:app-transfer-final}
\end{equation}
The second arrow is induced by
$U_\ell\xrightarrow{\sim}U^-$, $u\mapsto u^{-T}$.
\end{proposition}

\begin{proof}
The first isomorphism is the two-sided Grosshans transfer for the two basic
affine spaces, evaluated at the standard flags; see \cite{Grosshans}.
Substituting $a=u^{-T}$ identifies the resulting $U^-\times U$-action with
\eqref{eq:two-sided-action}.  The central scaling torus and the two terminal
diagonal tori commute with these maps, so the isomorphism preserves the full
multigrading.
\end{proof}

In particular, the determinant of the $r$-th central arrow maps to
$\Delta_r=\det A_r$.

\subsection{Evaluation of the Schofield presentations}

We use Schofield's determinantal construction of quiver semi-invariants
\cite{SchofieldSemi}; the particular presentations and paths below are those
of \cite{FeiKroneckerII}.  On the determinant-open locus, the latter paper
adjoins formal
inverses of the central arrows.  If $I=(r_1,\ldots,r_{2q+1})$ has odd length, the
corresponding path evaluates as
\[ A_I=A_{r_1}A_{r_2}^{-1}A_{r_3}\cdots A_{r_{2q}}^{-1}A_{r_{2q+1}}. \]
For $i+j=k$, the primal presentation used there is
\[ f_{i,j}^{I,J}:P_k\longrightarrow P_{-i}\oplus P_{-j}, \]
and the reflected presentation reverses the paths and interchanges source
and target.

\begin{lemma}[matrix form of the Schofield semi-invariants]
\label{lem:app-Schofield-matrix}
After restricting both flag arms to the standard flags, the Schofield
semi-invariant of $f_{i,j}^{I,J}$ is, up to a nonzero scalar,
\begin{equation}
 \det
 \begin{bmatrix}
  P_iA_IJ_k\\[1mm]
  P_jA_JJ_k
 \end{bmatrix}.
 \label{eq:app-left-block}
\end{equation}
The reflected semi-invariant is, up to the corresponding scalar,
\begin{equation}
 \det\left[
  P_kA_{\overleftarrow I}J_i
  \ \middle|\ 
  P_kA_{\overleftarrow J}J_j
 \right].
 \label{eq:app-right-block}
\end{equation}
Reflection of the flagged quiver exchanges these two determinants and
becomes matrix transposition under \eqref{eq:app-transfer-final}.
\end{lemma}

\begin{proof}
On the standard negative arm, the unique path from the $i$-dimensional
vertex to the central vertex evaluates as $J_i$; on the standard positive
arm, the unique path from the central vertex to the $i$-dimensional vertex
evaluates as $P_i$.  By definition, the Schofield matrix is obtained from
the transpose of the presentation matrix by evaluating each path.  This
gives the vertical block \eqref{eq:app-left-block}.  Reflection reverses
all paths and interchanges source and target, giving the horizontal block
\eqref{eq:app-right-block}.
\end{proof}

Table~\ref{tab:flagged-matrix-dictionary} records the dictionary.  We use
the notation of \cite{FeiKroneckerII}:
$({}_{i,j}^{n})$ for the primal triangle and
$({}_{i,j}^{n\vee})$ for the reflected triangle.

\begin{table}[H]
\centering
\small
\renewcommand{\arraystretch}{1.35}
\begin{tabular}{@{}L{.20\textwidth}|L{.72\textwidth}@{}}
\textbf{Index in \cite{FeiKroneckerII}} & \textbf{presentation and image after standard-flag evaluation}\\ \hline
$({}_{i,j}^{n})$, $k=i+j$
 & $f_{i,j}^{n,\mathbf q_n}:P_k\to P_{-i}\oplus P_{-j}$, and
   \(\displaystyle
   s(f_{i,j}^{n,\mathbf q_n})\longmapsto
   \det\!\begin{bmatrix}P_iA_nJ_k\\ P_jA_{\mathbf q_n}J_k\end{bmatrix}
   =\phi_{i,j}^{(n),\mathrm L,\circ}.\)\\[2mm]
$({}_{i,j}^{n\vee})$, $k=i+j$
 & The reflected presentation $\check f_{i,j}^{n,\mathbf q_n}$ maps to
   \(\displaystyle
   \det\!\left[P_kA_nJ_i\ \middle|\
                    P_kA_{\overleftarrow{\mathbf q_n}}J_j\right]
   =\phi_{i,j}^{(n),\mathrm R,\circ}.\)\\[2mm]
$n$
 & The coefficient $\det M(a_n)$ maps to $\Delta_n=\det A_n$.
\end{tabular}
\caption{Labels from \cite{FeiKroneckerII}, presentations, and matrix determinants.}
\label{tab:flagged-matrix-dictionary}
\end{table}

\subsection{Minimal lift, boundary identifications, and corner corrections}

Definition~2.2 and Lemma~2.5 of \cite{FeiKroneckerII} give the minimal
lift exactly, not only up to its divisor.  Under Table~\ref{tab:flagged-matrix-dictionary} it becomes
\begin{equation}
 \widetilde s_{i,j}^{n}\longmapsto
 \begin{cases}
  \phi_{i,0}^{(n),\mathrm L,\circ},&j=0,\\[1mm]
  \Pi_n\phi_{i,j}^{(n),\mathrm L,\circ},
    &j>0\text{ and }(i+j<\ell\text{ or }n\text{ even}),\\[1mm]
  \phi_{i,j}^{(n-1),\mathrm L},
    &j>0,\ i+j=\ell,\ n\text{ odd},
 \end{cases}
 \label{eq:app-minimal-lift-dictionary}
\end{equation}
with the transpose formula for $\widetilde{\check s}_{i,j}^{n}$.  Thus
\eqref{eq:app-minimal-lift-dictionary} is precisely
\eqref{eq:regular-left}--\eqref{eq:regular-right}.  The identifications of
edge labels are:
\begin{equation}
\begin{array}{c|c|c}
 \text{edge in \cite{FeiKroneckerII}}&\text{identification in \cite{FeiKroneckerII}}&\text{matrix notation}\\ \hline
 \text{common base}&({}_{i,0}^{n})=({}_{i,0}^{n\vee})
    &\det(P_iA_nJ_i)\\
 n\ge4\text{ even}&({}_{0,j}^{n})=({}_{0,j}^{n-1})
    &x_{0,j}^{(n)}=x_{0,j}^{(n-1)}\\
 n\ge3\text{ odd},\ i+j=\ell
    &({}_{i,j}^{n})=({}_{i,j}^{n-1})
    &x_{i,j}^{(n)}=x_{i,j}^{(n-1)}
\end{array}
\label{eq:app-edge-identification-table}
\end{equation}
These are exactly the boundary-identification conventions used in
\eqref{eq:even-identification}--\eqref{eq:odd-identification}.

The only lifted exchange relations requiring a new determinant arrow are
the three corner families in \cite[Corollaries~5.3--5.4]{FeiKroneckerII};
they are listed below.
\begin{equation}
\begin{array}{c|c|c}
 \text{corner in \cite{FeiKroneckerII}}&\text{coefficient arrow}&\text{matrix relation}\\ \hline
 ({}_{\ell-1,0}^{n})&({}_{\ell-1,0}^{n})\longrightarrow n
   &\Delta_n\text{ at the base corner}\\
 n\text{ even}:({}_{\ell-1,1}^{n}),({}_{\ell-1,1}^{n\vee})
   &n\longrightarrow\text{both}
   &c_{n;\ell-1,1}=\Delta_n\\
 n\ge3\text{ odd}:({}_{0,\ell-1}^{n}),({}_{0,\ell-1}^{n\vee})
   &n\longrightarrow\text{both}
   &\kappa_{n;\ell-1}=\Delta_n
\end{array}
\label{eq:app-corner-correction-table}
\end{equation}
In the first row, substituting
the lift cases from \eqref{eq:app-minimal-lift-dictionary} explains the
parity switch: for odd $n$ the boundary variable is inherited from
block $n-1$, so the normalized determinant coefficient occurs in the
opposite exchange monomial.  No other determinant correction is present.

Lemma~\ref{lem:app-Schofield-matrix} and
Table~\ref{tab:flagged-matrix-dictionary} identify the unlifted functions.
Equation~\eqref{eq:app-minimal-lift-dictionary} gives the minimal lifts, while
\eqref{eq:app-edge-identification-table} and
\eqref{eq:app-corner-correction-table} give the edge identifications and
coefficient arrows.  Thus the transfer isomorphism carries the full lifted
seed of \cite{FeiKroneckerII}, including all exchange relations, to the
matrix seed of Sections~\ref{sec:matrix-invariants} and \ref{sec:folding};
reflection of the flagged presentation becomes matrix transposition.

\section{Transverse curves for the remaining coprimality cases}
\label{app:exceptional-coprimality}

The three families in \eqref{eq:exceptional-coprimality-families} admit
curves transverse to both the initial divisor and its exchange numerator.
Matrices not specified below are identities; the stated blocks are extended by an
identity block, and $\doteq$ denotes equality up to a nonzero scalar.

For a mutable direction $u$, write
\[ N_u=M_{u,+}+M_{u,-}=z_uz_u'. \]

\begin{proposition}\label{prop:exceptional-specializations}
Every mutable variable in
\eqref{eq:exceptional-coprimality-families} admits a curve $\gamma:\mathbb A^1\to X_{\ell,m}$ satisfying
\[ z_u(t)=t\eta(t),\qquad N_u(t)=tq(t), \qquad \eta(0)q(0)\ne0. \]
\end{proposition}

\begin{proof}
The permutation matrices below are symmetric involutions.  They are
realized in an alternating word by taking one earlier factor equal to the
chosen involution and all other factors equal to the identity.

\smallskip
\noindent\emph{Base-edge variable for even $n$.}
Take
\[ C_n(t)=\begin{pmatrix}t&0\\0&1\end{pmatrix}\oplus I_{\ell-2}. \]
Then
\[ z_{1,0}^{(n)}=t,\qquad z_{2,0}^{(n)}\doteq t, \qquad z_{0,1}^{(n)}\doteq1, \qquad z_{1,1}^{(n)}=0. \]
The base-edge numerator has order one.  This includes $\ell=2$, where
$z_{2,0}^{(n)}=\Delta_n=t$.

\smallskip
\noindent\emph{Base-edge variable for odd $n$.}
Let
\[ G=\begin{pmatrix}0&1\\1&0\end{pmatrix}\oplus I_{\ell-2}, \qquad C_n(t)=\begin{pmatrix}t&1\\1&1\end{pmatrix}\oplus I_{\ell-2}, \qquad D=t-1. \]
and choose the preceding word so that $A_{\mathbf q_n}=GC_n$.  For
$\ell\ge3$,
\[ z_{1,0}^{(n)}=t, \qquad z_{2,0}^{(n)}\doteq D, \qquad z_{0,1}^{(n)}\doteq1, \qquad z_{1,1}^{(n)}\doteq D. \]
Thus $N_{1,0}=aD+bD^2$ with $a,b\ne0$.  Polynomial regularity gives
$t\mid N_{1,0}$, so $a=b$ and $N_{1,0}=aDt$.  For $\ell=2$, the same
argument applied to the corner numerator $aD+b$ gives $N_{1,0}=at$.
Interchanging $C_1$ and $C_2$ in the even $n=2$ curve treats
$z_{0,1}^{(2)}$.

\smallskip
\noindent\emph{Interior variable for even $n$.}
For $z_{1,2}^{(n)}$, take
\[
 C_n(t)=
 \begin{pmatrix}
  0&1&t&0\\
  1&0&0&1\\
  t&0&0&0\\
  0&1&0&0
 \end{pmatrix}\oplus I_{\ell-4}.
\]
The relevant chamber determinants are
\[ z_{1,2}^{(n)}\doteq t, \qquad (z_{1,1},z_{2,2},z_{0,3};z_{1,3},z_{0,2},z_{2,1}) \doteq(-1,t,1;0,1,0). \]
The two terms in \eqref{eq:interior} therefore have sum of order one.

\smallskip
\noindent\emph{Interior variable for odd $n$ when $\ell=4$.}
Let $G$ exchange $e_1$ with $e_3$ and $e_2$ with $e_4$.  Choose the preceding
matrices so that $C_{n-1}=I_4$ and $A_{\mathbf q_{n-1}}=G$; hence
$A_{\mathbf q_n}=GC_n$.  Put
\[
 C_n(t)=
 \begin{pmatrix}
  1&0&0&0\\
  0&0&1&0\\
  0&1&0&t\\
  0&0&t&1
 \end{pmatrix}.
\]
Then
\[ z_{1,2}^{(n)}\doteq t, \qquad (z_{1,1},z_{2,2},z_{0,3};z_{1,3},z_{0,2},z_{2,1}) \doteq(1,1,t;0,0,-1). \]
Thus the first term in \eqref{eq:interior} has order one and the second
vanishes.

\smallskip
\noindent\emph{Interior variable for odd $n$ when $\ell\ge5$.}
Let $G$ exchange $e_1$ with $e_4$ and $e_2$ with $e_5$,
and put
\[
 C_n(t)=
 \begin{pmatrix}
  0&-1&0&-1&0\\
  -1&1&0&0&1\\
  0&0&-1&0&t\\
  -1&0&0&0&0\\
  0&1&t&0&0
 \end{pmatrix}\oplus I_{\ell-5}.
\]
With $A_{\mathbf q_n}=GC_n$,
\[ z_{1,2}^{(n)}\doteq-t, \qquad (z_{1,1},z_{2,2},z_{0,3};z_{1,3},z_{0,2},z_{2,1}) \doteq(-1,t,1;1,-1,0). \]
Again the numerator in \eqref{eq:interior} has order one.
\end{proof}

\section{Row formulas for a folded half-diamond}
\label{app:triangular-mutation}

The diagonal-removal lemma, Lemma~\ref{lem:folded-half-diamond-removal},
rests on row formulas for a triangular mutation sequence.  Only the mutable
principal matrix enters the argument; determinant coefficients and frozen
boundary columns are omitted.

Let
\[ \mathsf T_N=\{(a,b)\in\mathbb Z_{\ge0}^2:1\le a+b\le N\}, \qquad \mathsf T_N^+=\{(a,b)\in\mathsf T_N:a>0\}. \]
For a label $(a,b)$, let $e_{a,b}$ be the corresponding unit row, and set
$e_{a,b}=0$ when the label lies outside $\mathsf T_N$.  Consider the
extended matrix with mutable rows indexed by $\mathsf T_N^+$ whose rows are
\begin{align}
 r_{a,0}
 &=e_{a+1,0}+2e_{a-1,1}-e_{a-1,0}-2e_{a,1},
 \label{eq:appendix-base-row}\\
 r_{a,b}
 &=e_{a,b-1}+e_{a+1,b}+e_{a-1,b+1}
   -e_{a,b+1}-e_{a-1,b}-e_{a+1,b-1}
 \quad(b>0).
 \label{eq:appendix-interior-row}
\end{align}
These are the mutable rows of an exposed even folded half-diamond.
Replacing the entire matrix by its negative only reverses every source
and sink, so the chosen orientation causes no loss of generality.

\begin{figure}[H]
\centering
\begin{tikzpicture}[
  every node/.style={font=\scriptsize},
  mutated/.style={circle,draw,minimum size=5.5mm,inner sep=0pt},
  plain/.style={circle,draw,minimum size=4.5mm,inner sep=0pt},
  edge/.style={line width=.35pt}
]
  \node[mutated] (v30) at (0,3) {$1$};
  \node[mutated] (v31) at (1.15,3) {$2$};
  \node[mutated] (v32) at (2.30,3) {$3$};
  \node[mutated] (v33) at (3.45,3) {$4$};
  \node[mutated] (v20) at (.575,2) {$5$};
  \node[mutated] (v21) at (1.725,2) {$6$};
  \node[mutated] (v22) at (2.875,2) {$7$};
  \node[mutated] (v10) at (1.15,1) {$8$};
  \node[mutated] (v11) at (2.30,1) {$9$};
  \node[plain,line width=.8pt] (d) at (1.725,0) {$d_i$};
  \node[plain] (p) at (.05,1) {};
  \node[plain] (c) at (-.525,2) {};
  \node[plain] (t) at (-1.10,3) {};

  \foreach \x/\y in {v30/v31,v31/v32,v32/v33,
                       v20/v21,v21/v22,v10/v11,
                       v30/v20,v31/v20,v31/v21,v32/v21,
                       v32/v22,v33/v22,v20/v10,v21/v10,
                       v21/v11,v22/v11,v10/d,v11/d,
                       t/c,c/p,p/d}
    \draw[edge] (\x)--(\y);
  \draw[dashed] (t)--(p);
  \node[anchor=east] at (-1.20,2) {adjacent side};
  \draw[-{Latex[length=2mm]},line width=.6pt]
    (3.9,3.25) .. controls (4.45,2.25) and (3.85,.95) .. (2.65,.25);
  \node[anchor=west,align=left] at (4.0,1.85)
    {anti-diagonals\\from top to bottom};
  \node[below=2pt of d] {$[r;i,0]$};
\end{tikzpicture}
\caption{The anti-diagonal mutation order when $N-i=3$.  The numbered
vertices are $v_{q,u}=(i+u,q-u)$ in mutation order; the numbers do not
indicate arrow orientation.}
\label{fig:half-diamond-mutation-order}
\end{figure}

Fix $1\le i<N$, delete the base vertices
$(1,0),\ldots,(i-1,0)$, and put
\[ h:=N-i, \qquad c:=(i-1,h+1), \]
\[ v_{q,u}:=(i+u,q-u) \qquad(1\le q\le h,\ 0\le u\le q). \]
The mutation sequence \eqref{eq:positive-removal-sequence}, with the block
label suppressed, is
\[ (v_{h,0},v_{h,1},\ldots,v_{h,h}; v_{h-1,0},\ldots,v_{h-1,h-1};\ldots; v_{1,0},v_{1,1}). \]

\begin{lemma}\label{lem:anti-diagonal-row-formula}
Immediately before mutating $v_{q,u}$ in the preceding sequence, its complete
row is the following.  In the middle lines below, put $a=i+u$ and
$b=q-u$.  For the top anti-diagonal $q=h$,
\begin{equation}
 r_{v_{h,u}}=
 \begin{cases}
  -e_{i-1,h}+e_c+e_{i,h-1}-e_{i+1,h-1},&u=0,\\
  e_c-e_{a-1,b+1}+e_{a,b-1}-e_{a+1,b-1},&0<u<h,\\
  2e_c+e_{N-1,0}-2e_{N-1,1},&u=h.
 \end{cases}
 \label{eq:top-mutation-rows}
\end{equation}
For $1\le q<h$,
\begin{equation}
 r_{v_{q,u}}=
 \begin{cases}
  -e_{i-1,q}+e_{i,q-1}+e_{i,q+1}-e_{i+1,q-1},&u=0,\\
  -e_{a-1,b+1}+e_{a,b-1}+e_{a,b+1}-e_{a+1,b-1},&0<u<q,\\
  e_{i+q-1,0}-2e_{i+q-1,1}+e_{i+q+1,0},&u=q.
 \end{cases}
 \label{eq:lower-mutation-rows}
\end{equation}
All other entries are zero.
\end{lemma}

\begin{proof}
We use the extended-matrix mutation formula
\begin{equation}
 b'_{xj}=b_{xj}+[b_{xv}]_+[b_{vj}]_+
                 -[-b_{xv}]_+[-b_{vj}]_+
 \qquad(x,j\ne v),
 \label{eq:appendix-mutation-formula}
\end{equation}
together with $b'_{xv}=-b_{xv}$.  We prove the assertion in the stated
mutation order and keep the \emph{complete} rows, including columns on the
adjacent unmutated side.

The first row in \eqref{eq:top-mutation-rows} is
\eqref{eq:appendix-interior-row} for $(a,b)=(i,h)$: the two terms outside
the triangle vanish because $i+h=N$.  For the vertex $v=(a,b)=v_{q,u}$ being mutated, with $b>0$, set
\[ R:=(a+1,b-1),\qquad D:=(a,b-1). \]
Among vertices occurring later in the sequence, the row of $v$ meets only
$R$
and $D$.  Consequently a mutation at $v$ changes no other future row.  It
therefore suffices to verify the two transfers to $R$ and $D$, together
with the valued terminal transfer when $R$ is a base vertex.

\smallskip
\noindent\emph{Transfer to the right.}
Assume first that $b>1$.  Immediately before $\mu_v$, the row of
$R=(a+1,b-1)$ is
\begin{equation}
 r_R=
 \begin{cases}
  e_v-e_{a,b-1}+e_{a+1,b-2}-e_{a+2,b-2},&q=h,\\
  e_v-e_{a,b-1}-e_{a,b+1}
      +e_{a+1,b-2}+e_{a+1,b}-e_{a+2,b-2},&q<h.
 \end{cases}
 \label{eq:right-intermediate-row}
\end{equation}
For $q=h$ this is the original row
\eqref{eq:appendix-interior-row}; for $q<h$ it is the intermediate row obtained after mutating
the vertex $(a+1,b)$ directly above $R$, as recorded in the
downward transfer below.  Since $b_{Rv}=1$, substitution of
\eqref{eq:top-mutation-rows} or \eqref{eq:lower-mutation-rows} into
\eqref{eq:appendix-mutation-formula} cancels the negative terms involving
$a,b$ in \eqref{eq:right-intermediate-row}, changes $e_v$ to $-e_v$, and gives
precisely the next middle row in \eqref{eq:top-mutation-rows} or
\eqref{eq:lower-mutation-rows}.  Thus the induction advances from
$v_{q,u}$ to $v_{q,u+1}$.

If $b=1$, then $R=(a+1,0)$ is a transpose-fixed base vertex and
$b_{Rv}=2$.  Before $\mu_v$ its row is
\[
 r_R=
 \begin{cases}
  -e_{a,0}+2e_v,&q=h,\\
  -e_{a,0}+2e_v-2e_{a,2}+e_{a+2,0},&q<h.
 \end{cases}
\]
Formula \eqref{eq:appendix-mutation-formula} adds twice the positive part
of $r_v$ and reverses the $v$-entry.  It gives respectively
\[ 2e_c+e_{a,0}-2e_v, \qquad e_{a,0}-2e_v+e_{a+2,0}, \]
which are the terminal rows in \eqref{eq:top-mutation-rows} and
\eqref{eq:lower-mutation-rows}.

\smallskip
\noindent\emph{Transfer downward.}
Suppose $b>1$.  If $u=0$, a direct substitution in
\eqref{eq:appendix-mutation-formula} gives
\[ r_D'=-e_{i-1,q-1}+e_{i,q-2}+e_v-e_{i+1,q-2}, \]
which is the first row of \eqref{eq:lower-mutation-rows} for the next
anti-diagonal.  If $u>0$, the same calculation gives the intermediate row
\begin{equation}
 \overline r_D=
 -e_{a-1,b-1}+e_{a-1,b}-e_{a-1,b+1}
 +e_{a,b-2}+e_{a,b}-e_{a+1,b-2}.
 \label{eq:intermediate-lower-row}
\end{equation}
Before $D$ is reached in the next anti-diagonal, its left predecessor
$P=(a-1,b)$ is mutated.  The only positive entries of $r_P$ that meet
\eqref{eq:intermediate-lower-row} are
$e_{a-1,b-1}+e_{a-1,b+1}$, while
$b_{DP}=1$.  Mutation at $P$ cancels the two negative terms in
\eqref{eq:intermediate-lower-row}, reverses the $P$-entry, and leaves
\[ -e_{a-1,b}+e_{a,b-2}+e_{a,b}-e_{a+1,b-2}, \]
which is the required middle row of \eqref{eq:lower-mutation-rows} for $D$.
The second line of \eqref{eq:right-intermediate-row} follows because
$R$ is the downward successor of the vertex $(a+1,b)$ in the preceding
anti-diagonal.

For the downward successor of the terminal pair, let
$v=(a,1)$ and $R=(a+1,0)$.  After the consecutive mutations
$\mu_v\mu_R$, the row of $D=(a,0)$ is
\begin{equation}
 \overline r_D=
 -e_{a-1,0}+2e_{a-1,1}-2e_{a-1,2}+e_{a+1,0}.
 \label{eq:intermediate-lower-base-row}
\end{equation}
When the next anti-diagonal is processed, its last interior vertex
$P=(a-1,1)$ is mutated immediately before $D$.  Here $b_{DP}=2$, and the
positive part of $r_P$ relevant to
\eqref{eq:intermediate-lower-base-row} is
$e_{a-1,0}+e_{a-1,2}$.  Hence mutation at $P$ turns
\eqref{eq:intermediate-lower-base-row} into
$e_{a-1,0}-2e_{a-1,1}+e_{a+1,0}$,
the terminal row in \eqref{eq:lower-mutation-rows}.

The right and downward transfers account for every later row changed by
each mutation.  Starting from $v_{h,0}$, these calculations determine the
complete rows on the top anti-diagonal and then those on each lower
anti-diagonal.  This proves
\eqref{eq:top-mutation-rows}--\eqref{eq:lower-mutation-rows} by induction.
\end{proof}

\begin{proposition}\label{prop:triangular-mutation-sequence}
After deleting $(1,0),\ldots,(i-1,0)$ and applying
$\mathbf p_i^+$, the row of $(i,0)$ is
\begin{equation}
 r_{i,0}=e_{i+1,0}
 \qquad(1\le i<N).
 \label{eq:appendix-final-base-row}
\end{equation}
For $i=N$, the sequence is empty and, after the previous base vertices have
been deleted,
\begin{equation}
 r_{N,0}=2e_{N-1,1}.
 \label{eq:appendix-terminal-base-row}
\end{equation}
The same conclusions hold after deleting any additional unmutated boundary
vertices.  After deleting $(i,0)$, applying the mutations in reverse order
restores the original matrix on all surviving vertices.
\end{proposition}

\begin{proof}
Assume $i<N$.  Before the last anti-diagonal is reached, the row of
$d=(i,0)$ is unchanged, because none of the vertices with $q>1$ is adjacent
to it.  Since $(i-1,0)$ has already been deleted,
\begin{equation}
 r_d=2e_{i-1,1}-2e_{i,1}+e_{i+1,0}.
 \label{eq:appendix-tracked-row}
\end{equation}
Put $x=v_{1,0}=(i,1)$ and $y=v_{1,1}=(i+1,0)$.
Lemma~\ref{lem:anti-diagonal-row-formula} gives
\[ r_x=-e_{i-1,1}+e_d+t-e_y, \]
where $t=e_{i-1,2}$ if $h=1$ and $t=e_{i,2}$ if $h>1$; the term $t$ is
absent when its label is absent.  Since $b_{dx}=-2$, mutation at $x$
cancels the $2e_{i-1,1}$ term in
\eqref{eq:appendix-tracked-row}, changes the sign of the $x$-entry, and
gives
$r_d'=2e_x-e_y$.
Immediately before its mutation, the terminal row has the form
\[ r_y=e_d-2e_x+s, \]
where $s=2e_c$ if $h=1$ and $s=e_{i+2,0}$ if $h>1$.  In either case all
entries of $s$ are nonnegative.  Since $b_{dy}=-1$, mutation at $y$ cancels
$2e_x$ and reverses the $y$-entry.  Thus the row of $d$ passes through the
following three values:
\[ \underbrace{2e_{i-1,1}-2e_x+e_y}_{\text{before the last anti-diagonal}} \xrightarrow{\ \mu_x\ } \underbrace{2e_x-e_y}_{\text{after $\mu_x$}} \xrightarrow{\ \mu_y\ } \underbrace{e_y}_{\text{source row}}. \]
which proves \eqref{eq:appendix-final-base-row}.

For $i=N$, formula \eqref{eq:appendix-base-row} has only the term
$2e_{N-1,1}$ after $(N-1,0)$ is deleted, which proves
\eqref{eq:appendix-terminal-base-row}.  Principal restriction commutes with
mutation away from the removed vertex, so any additional unmutated
boundary deletion restricts these identities.  The same commutation
shows that, after $d$ is deleted, the reverse sequence is the inverse of the
forward sequence on the restricted matrix.
\end{proof}


\section{Folding an even boundary path}
\label{app:even-boundary-path}

For Proposition~\ref{prop:optimized-boundary-seeds}, retain the mutable
principal matrix and the orbit of the frozen column being optimized.  The
other frozen columns do not occur in the mutation formulas below.

\subsection{Orbit mutation with one frozen orbit}

Let $\widetilde B$ be a $\tau$-equivariant mutable principal matrix, let
$\bar f$ be a frozen orbit, and retain the columns indexed by the members of
$\bar f$.  Its folded column is
\[ c_{\bar x,\bar f}:=\sum_{f'\in\bar f}\widetilde b_{x f'}, \qquad x\in\bar x. \]
For a set $S$ of mutable vertices, set
\[ \mathcal N_{\bar f}(S) :=S\cup\left\{x:\ \widetilde b_{xs}\ne0\text{ for some }s\in S \text{ or } \widetilde b_{x f'}\ne0\text{ for some }f'\in\bar f\right\}. \]

\begin{lemma}\label{lem:frozen-column-locality}
Under a mutation sequence supported in $S$, every mutable vertex outside
$\mathcal N_{\bar f}(S)$ remains nonadjacent to the unmutated members of
$S$ and retains zero entries in the columns of $\bar f$.  Hence the
principal matrix and the folded $\bar f$-column are determined by their
restriction to $\mathcal N_{\bar f}(S)$ together with the columns of $\bar f$.
\end{lemma}

\begin{proof}
Initially an outside row has zero entries both on $S$ and in the columns
of $\bar f$.  In the mutation formula, a correction to either kind of entry
contains the row's entry at the vertex being mutated, which is zero.
Induction on the mutation sequence proves the claim.
\end{proof}

\begin{lemma}\label{lem:frozen-column-orbit-mutation}
Let $\bar k$ be a mutable orbit.  Suppose, for a representative $k$, that
\begin{align}
 \widetilde b_{k,\tau k}&=0,
 \label{eq:frozen-fold-no-orbit-arrow}\\
 \widetilde b_{xk}\widetilde b_{x,\tau k}&\ge0
 &&\text{for every mutable row }x,
 \label{eq:frozen-fold-row-sign}\\
 \widetilde b_{k f}\widetilde b_{k,\tau f}&\ge0
 &&\text{for a representative }f\in\bar f.
 \label{eq:frozen-fold-column-sign}
\end{align}
Then simultaneous mutation at the members of $\bar k$ folds correctly on
the mutable principal matrix and on the orbit-sum column of $\bar f$.
\end{lemma}

\begin{proof}
By Lemma~\ref{lem:orbit-mutation-folding}, the first two conditions fold
the principal matrix.  Equivariance identifies the four entries between
$\bar k$ and $\bar f$ as
\[ \begin{pmatrix}a&b\\ b&a\end{pmatrix} \]
after choosing representatives (with the evident degeneration for a
singleton orbit).  Condition \eqref{eq:frozen-fold-column-sign} says that
$a$ and $b$ have one sign.  Together with
\eqref{eq:frozen-fold-row-sign}, this allows positive and negative parts to
pass through both orbit sums.  Writing
\(
 \Phi(r,s)=[r]_+[s]_+-[-r]_+[-s]_+
\), one obtains, for $\bar x\ne\bar k$,
\[ c'_{\bar x,\bar f} =c_{\bar x,\bar f} +\sum_{k'\in\bar k}\sum_{f'\in\bar f} \Phi(\widetilde b_{xk'},\widetilde b_{k'f'}) =c_{\bar x,\bar f} +\Phi(b_{\bar x,\bar k},c_{\bar k,\bar f}). \]
The entry at $\bar k$ changes sign, which is the mutation rule for the
folded $\bar f$-column.
\end{proof}

\subsection{Local sign conditions along a boundary path}

The boundary paths of \cite{FeiKroneckerII} are defined after twisting the
two open triangles of each odd diamond.  Pair each twist sequence, orbit by
orbit, with its transpose, and process the odd diamonds in the order used
there.

\begin{lemma}\label{lem:folded-preparatory-twists}
For any frozen orbit $\bar f$, the paired preparatory twists fold
correctly on the mutable principal matrix and the $\bar f$-column.  Their
result is the corresponding restriction of the fold of the twisted disk
seed in \cite{FeiKroneckerII}.
\end{lemma}

\begin{proof}
The two open triangles of one odd diamond are transpose images, meet only
along the fixed diagonal, and have no arrows between their open interiors.
Before each paired mutation, a mutable row is incident with at most one
member of the pair, unless it lies on the fixed diagonal, in which case the
two incidences are equal.  The frozen orbit is incident with one side, the
transpose side, both with the same sign, or neither.  Conditions
\eqref{eq:frozen-fold-no-orbit-arrow}--
\eqref{eq:frozen-fold-column-sign} therefore hold and are preserved by the
separated-support argument of Lemma~\ref{lem:separated-pair}.  Distinct
odd-diamond interiors are disjoint, so the twists may be processed
successively.
\end{proof}

In the twisted disk seed, put
\[ f:=\bar v, \qquad f^{\star}:=\overline{v^{\star}}. \]
Orient the full boundary path $p_{v^{\star}}$ from $v^{\star}$ toward the
terminal vertex $v$.  Pair its mutation sequence orbit by orbit with the
transpose sequence, and let $S$ be the set of central vertices that occur.
At each cross, mutate the center first; on each remaining maximal hive strip,
use the mutation sequence of \cite[Lemma~6.15]{FeiKroneckerII}.

A pair whose product is nonnegative has, up to a common sign, one of
the forms
\[ (d\epsilon,0),\qquad(0,d\epsilon),\qquad (d_1\epsilon,d_2\epsilon),\qquad(0,0), \qquad d_1,d_2>0. \]
The row-sign condition is local.  It suffices to examine the minimal
extended subquivers $\rho_s$ of \cite[Lemmas~6.15--6.16]{FeiKroneckerII},
which contain every arrow incident with the central vertex being mutated
and all later central vertices.  A row absent from the local forms below is
zero in the two mutation columns.

\begin{lemma}\label{lem:straight-strip-local-matrix}
Let $\{k,\tau k\}$ be the nonfixed orbit being mutated in a
$\tau$-stable straight strip.  Up to simultaneous reversal of all
arrows, every nonzero incidence with the orbit is one of the
following oriented patterns:
\begin{equation*}
\begin{array}{c@{\qquad}c@{\qquad}c}
 x\longrightarrow k,
 &x\longrightarrow k\quad\text{and}\quad x\longrightarrow\tau k,
 &x\longrightarrow k,\quad \tau x\longrightarrow\tau k,\\[-1mm]
 \text{one mutation column}
 &\text{both mutation columns}
 &\text{transpose-paired rows}.
\end{array}
\end{equation*}
After allowing the positive valuations created by orbit summation, the
complete block on a nonfixed row orbit $\{x,\tau x\}$ is therefore
\begin{equation}
 \epsilon
 \begin{pmatrix}d&0\\0&d\end{pmatrix}
 \quad\text{or}\quad
 \epsilon
 \begin{pmatrix}d&d'\\d'&d\end{pmatrix},
 \qquad
 \epsilon\in\{1,-1\},\quad d>0,\quad d'\ge0,
 \label{eq:straight-strip-orbit-blocks}
\end{equation}
and a transpose-fixed row is $\epsilon(d,d')$ with $d,d'>0$.  A
row absent from these blocks is zero in both mutation columns.  The same
oriented normal forms hold for the two columns of the chosen frozen orbit.
In particular,
\begin{equation}
 \widetilde b_{xk}\widetilde b_{x,\tau k}\ge0,
 \qquad
 \widetilde b_{k f}\widetilde b_{k,\tau f}\ge0,
 \qquad
 \widetilde b_{k,\tau k}=0.
 \label{eq:straight-strip-sign-conclusion}
\end{equation}
\end{lemma}

\begin{proof}
The minimal extended subquiver $\rho_s$ in the proof of
\cite[Lemma~6.15]{FeiKroneckerII} is a chain of oriented rhombi.  Choose the
orientation in which the side rows of the rhombus under consideration point
toward its central vertex; reversing the disk reverses all signs.  A side row,
a terminal row, or a later central row lying strictly in the chosen open
half is incident only with $k$, and its transpose is incident only with
$\tau k$.  This gives the diagonal block in
\eqref{eq:straight-strip-orbit-blocks}.  A row on the fold axis, or the row
obtained by identifying the two boundary side rows, is a side of both mirror
rhombi.  The cyclic orientations agree along this row, so both arrows point
toward the orbit being mutated, or both point away after simultaneous
reversal.
This gives the second block in \eqref{eq:straight-strip-orbit-blocks}.

Orbit summation may multiply entries by positive integers, but it does not
change their signs.  The defining property of $\rho_s$ accounts for every
arrow incident with the central vertex being mutated and the later central
vertices.  The strip mutation sequence replaces $\rho_s$ by the next
oriented rhombus chain of the same form, so
the stated alternatives hold before every mutation in the strip.  The
chosen terminal row satisfies the same alternatives.  Finally, $k$ and
$\tau k$ lie in distinct mirror rhombi and are nonadjacent.
\end{proof}

At a cross, use the notation of
\cite[Lemma~6.16]{FeiKroneckerII}:
\[ c=i,\qquad h_-=i-1,\quad h_+=i+1,\qquad v_-=(i-1)',\quad v_+=(i+1)'. \]
Thus $h_-,c,h_+$ lie on the unprimed strip and $v_-,c,v_+$ on the primed
strip.  Up to simultaneous reversal of all arrows, mutation at $c=i$ has
the following local form.
\begin{center}
\begin{tikzpicture}[
  x=1.05cm,y=1.05cm,
  vertex/.style={draw,fill=white,rounded corners=2.5pt,minimum width=8.5mm,
                 minimum height=6mm,inner sep=.8pt,font=\small},
  arr/.style={-{Latex[length=1.6mm]},line width=.55pt}]
\node[vertex] (c1) at (0,0) {$i$};
\node[vertex] (hm1) at (-1.55,0) {$i-1$};
\node[vertex] (hp1) at (1.55,0) {$i+1$};
\node[vertex] (vm1) at (0,-1.45) {$(i-1)'$};
\node[vertex] (vp1) at (0,1.45) {$(i+1)'$};
\draw[arr] (hm1)--(c1);
\draw[arr] (c1)--(hp1);
\draw[arr] (vm1)--(c1);
\draw[arr] (c1)--(vp1);
\draw[arr] (hp1)--(vm1);
\draw[arr] (vp1)--(hm1);

\node at (2.55,0) {$\xrightarrow{\ \mu_i\ }$};

\node[vertex] (c2) at (5.1,0) {$i$};
\node[vertex] (hm2) at (3.55,0) {$i-1$};
\node[vertex] (hp2) at (6.65,0) {$i+1$};
\node[vertex] (vm2) at (5.1,-1.45) {$(i-1)'$};
\node[vertex] (vp2) at (5.1,1.45) {$(i+1)'$};
\draw[arr] (c2)--(hm2);
\draw[arr] (hp2)--(c2);
\draw[arr] (c2)--(vm2);
\draw[arr] (vp2)--(c2);
\draw[arr,bend left=24] (hm2) to (hp2);
\draw[arr,bend right=24,preaction={draw=white,line width=2.1pt}] (vm2) to (vp2);
\end{tikzpicture}
\end{center}
The opposite orientation is obtained by reversing every arrow.  In the
right-hand quiver, the only arrows among the four vertices adjacent to $i$
are $h_-\to h_+$ and $v_-\to v_+$.

\begin{lemma}[mutation at a cross]
\label{lem:cross-local-matrix}
After mutation at $c$, there are no arrows joining the unprimed and primed
adjacent pairs:
\begin{equation}
 \widetilde B_{\{h_-,h_+\},\{v_-,v_+\}}=0.
 \label{eq:cross-zero-block}
\end{equation}
For every later orbit mutated on either strip, all side rows,
transpose-fixed rows, and the chosen frozen orbit have one of the forms in
\eqref{eq:straight-strip-orbit-blocks}; rows on the other strip are zero in
the two mutation columns.  Hence the three conditions in
\eqref{eq:straight-strip-sign-conclusion} hold throughout both strip
mutation sequences.
\end{lemma}

\begin{proof}
Mutation at $c$ reverses the four incident arrows.  The paths
$h_-\to c\to h_+$ and $v_-\to c\to v_+$ create the two arrows along the
strips.  The mixed paths create arrows opposite to $v_+\to h_-$ and
$h_+\to v_-$, so those pairs cancel.  This proves
\eqref{eq:cross-zero-block}.

Each external side row belongs to one of the four oriented triangles
meeting the cross.  Relative to a later orbit on one strip, its row has one
nonzero entry in the two mutation columns.  A transpose-fixed side row, or a boundary
row identified with its transpose, meets the two mirror triangles with the
same orientation and therefore has entries of the same sign in the two
mutation columns.  These are the forms in
\eqref{eq:straight-strip-orbit-blocks}.  The chosen terminal row belongs to
one of the same triangles.  The minimal extended subquiver in
\cite[Lemma~6.16]{FeiKroneckerII} accounts for all incidences with the strip
mutation sequences.  Orbit summation changes entries only by positive factors, and
\eqref{eq:cross-zero-block} makes every central row on the other strip zero
in the mutation columns.
\end{proof}

\begin{lemma}[sign condition along a boundary path]
\label{lem:boundary-path-sign-invariant}
Immediately before every nonfixed orbit $\{k,\tau k\}$ in the
$\tau$-stable boundary mutation sequence:
\begin{enumerate}
 \item $k$ and $\tau k$ are not adjacent;
 \item $\widetilde b_{xk}\widetilde b_{x,\tau k}\ge0$ for every
       mutable row $x\in\mathcal N_{\bar f}(S)$;
 \item $\widetilde b_{k f}\widetilde b_{k,\tau f}\ge0$ for a
       representative $f\in\bar f$.
\end{enumerate}
Consequently every orbit mutation folds on the principal matrix and the
chosen frozen column.
\end{lemma}

\begin{proof}
On a maximal straight strip, Lemma~\ref{lem:straight-strip-local-matrix}
applies before every mutation in the strip mutation sequence.  At a cross, mutate the
center first; Lemma~\ref{lem:cross-local-matrix} then applies to both
remaining strips, and the zero cross-block permits the two strip mutation
sequences to be interleaved.  The local forms exhaust $\mathcal N_{\bar f}(S)$, while
rows outside this set remain zero by
Lemma~\ref{lem:frozen-column-locality}.  Conditions
\eqref{eq:frozen-fold-no-orbit-arrow}--
\eqref{eq:frozen-fold-column-sign} therefore hold before every nonfixed
orbit mutation.  Apply Lemma~\ref{lem:frozen-column-orbit-mutation}.  Fixed
orbits are singletons and require no folding check.
\end{proof}

\subsection{The full even boundary path}

\begin{proposition}[folding the even boundary path]
\label{prop:folded-even-boundary-path}
Assume $m$ is even and $v\ne v^{\star}$.  Then
$f=\bar v$ and $f^{\star}=\overline{v^{\star}}$ are distinct.  After the
preparatory twists, the full boundary-path mutation sequence of
\cite{FeiKroneckerII} on $p_{v^{\star}}$, paired with its transpose, folds
on the mutable principal matrix and the chosen frozen column.  For some
mutable index $u$ and positive integers $c,c'$, the terminal incidences are
\begin{equation}
 f^{\star}\xrightarrow{\ c'\ }u\xrightarrow{\ c\ }f.
 \label{eq:appendix-even-boundary-terminal}
\end{equation}
In particular, the folded $f$-column is $c e_u$.
\end{proposition}

\begin{proof}
Lemma~\ref{lem:folded-preparatory-twists} places the relevant principal
matrix and chosen frozen column in the fold of the twisted disk chart in \cite{FeiKroneckerII}.
Lemma~\ref{lem:boundary-path-sign-invariant} then folds every mutation in
the full path mutation sequence.  Because $m$ is even, $v^{\star}$ lies in
the same half of the terminal disk as $v$, while transposition sends it to $\check v=\tau v$ in the reflected half.
Thus $v\ne v^{\star}$ implies $f\ne f^{\star}$, and the other terminal
vertex of $p_{v^{\star}}$ is $v$.  The strip and cross calculations in
\cite[Lemmas~6.15, 6.16, and 6.18]{FeiKroneckerII} leave the three-vertex
path
\[ v^{\star}\longrightarrow\widetilde u\longrightarrow v \]
(up to simultaneous reversal when the path orientation is reversed).  We
orient the sequence toward $v$, so the final arrow points into $v$.  Since the terminal
orbits are distinct, orbit summing gives
\eqref{eq:appendix-even-boundary-terminal}.  The corresponding minimal extended
subquiver contains every mutable incidence with the terminal vertices, and
Lemma~\ref{lem:frozen-column-locality} excludes any omitted row.
\end{proof}


\bibliographystyle{amsalpha}
\bibliography{bosonic_plethysm}


%\begin{thebibliography}{99}
%
%\bibitem{BFZ}
%A.~Berenstein, S.~Fomin, and A.~Zelevinsky,
%\emph{Cluster algebras III: upper bounds and double Bruhat cells},
%Duke Math. J. \textbf{126} (2005), no.~1, 1--52.
%
%\bibitem{CaoString}
%P.~Cao,
%\emph{On exchange matrices from string diagrams},
%arXiv:2203.07822v2, 2022.
%
%\bibitem{CaoLi}
%P.~Cao and F.~Li,
%\emph{Uniformly column sign-coherence and the existence of maximal green
%sequences},
%J. Algebraic Combin. \textbf{50} (2019), no.~4, 403--417.
%
%\bibitem{CMMM}
%M.-W.~Cheung, T.~Magee, T.~Mandel, and G.~Muller,
%\emph{Valuative independence and cluster theta reciprocity},
%arXiv:2505.09585v2, 2025, revised 2026.
%
%\bibitem{DerksenWeyman}
%H.~Derksen and J.~Weyman,
%\emph{Semi-invariants of quivers and saturation for
%Littlewood--Richardson coefficients},
%J. Amer. Math. Soc. \textbf{13} (2000), no.~3, 467--479.
%
%\bibitem{DupontFolding}
%G.~Dupont,
%\emph{An approach to non-simply-laced cluster algebras},
%J. Algebra \textbf{320} (2008), no.~4, 1626--1661.
%
%\bibitem{Fsemi}
%J.~Fei,
%\emph{Cluster algebras and semi-invariant rings I. Triple flags},
%Proc. Lond. Math. Soc. (3) \textbf{115} (2017), no.~1, 1--32.
%
%\bibitem{FeiKroneckerII}
%J.~Fei,
%\emph{Cluster algebras, invariant theory, and Kronecker coefficients II},
%Adv. Math. \textbf{341} (2019), 536--582.
%
%\bibitem{FeiWeymanExtension}
%J.~Fei and J.~Weyman,
%\emph{Extending upper cluster algebras},
%arXiv:1707.04661, 2017.
%
%\bibitem{FFP}
%J.~Fei,
%\emph{Cluster algebras, invariant theory, and fermionic Plethysm},
%to appear.
%
%\bibitem{FockGoncharov}
%V.~V. Fock and A.~B. Goncharov,
%\emph{Moduli spaces of local systems and higher Teichm\"uller theory},
%Publ. Math. Inst. Hautes \'Etudes Sci. No.~103 (2006), 1--211.
%
%\bibitem{FockGoncharovEnsembles}
%V.~V. Fock and A.~B. Goncharov,
%\emph{Cluster ensembles, quantization and the dilogarithm},
%Ann. Sci. \'Ec. Norm. Sup\'er. (4) \textbf{42} (2009), no.~6, 865--930.
%
%\bibitem{FZ1}
%S.~Fomin and A.~Zelevinsky,
%\emph{Cluster algebras I: Foundations},
%J. Amer. Math. Soc. \textbf{15} (2002), no.~2, 497--529.
%
%\bibitem{FZ4}
%S.~Fomin and A.~Zelevinsky,
%\emph{Cluster algebras IV: Coefficients},
%Compos. Math. \textbf{143} (2007), no.~1, 112--164.
%
%\bibitem{GSV}
%M.~Gekhtman, M.~Shapiro, and A.~Vainshtein,
%\emph{Drinfeld double of $GL_n$ and generalized cluster structures},
%Proc. Lond. Math. Soc. (3) \textbf{116} (2018), no.~3, 429--484.
%
%\bibitem{GHKK}
%M.~Gross, P.~Hacking, S.~Keel, and M.~Kontsevich,
%\emph{Canonical bases for cluster algebras},
%J. Amer. Math. Soc. \textbf{31} (2018), no.~2, 497--608.
%
%\bibitem{Grosshans}
%F.~D. Grosshans,
%\emph{Algebraic Homogeneous Spaces and Invariant Theory},
%Lecture Notes in Math. 1673, Springer, Berlin, 1997.
%
%\bibitem{Macdonald}
%I.~G. Macdonald,
%\emph{Symmetric Functions and Hall Polynomials}, second ed.,
%Oxford Mathematical Monographs, Oxford University Press, Oxford, 1995.
%
%\bibitem{PopovVinberg}
%V.~L. Popov and E.~B. Vinberg,
%\emph{Invariant theory},
%in Algebraic Geometry IV, Encyclopaedia Math. Sci. 55,
%Springer, Berlin, 1994, 123--278.
%
%\bibitem{Rosenlicht}
%M.~Rosenlicht,
%\emph{Some basic theorems on algebraic groups},
%Amer. J. Math. \textbf{78} (1956), no.~2, 401--443.
%
%\bibitem{SchofieldSemi}
%A.~Schofield,
%\emph{Semi-invariants of quivers},
%J. London Math. Soc. (2) \textbf{43} (1991), no.~3, 385--395.
%
%\bibitem{YeSymmetric}
%J.~Y.~Ye,
%\emph{Cluster structures on $SL_n/SO_n$},
%arXiv:2607.14634, 2026.
%
%\end{thebibliography}

\end{document}
